%%
% Copyright (c) 2017 - 2025, Pascal Wagler;
% Copyright (c) 2014 - 2025, John MacFarlane
%
% All rights reserved.
%
% Redistribution and use in source and binary forms, with or without
% modification, are permitted provided that the following conditions
% are met:
%
% - Redistributions of source code must retain the above copyright
% notice, this list of conditions and the following disclaimer.
%
% - Redistributions in binary form must reproduce the above copyright
% notice, this list of conditions and the following disclaimer in the
% documentation and/or other materials provided with the distribution.
%
% - Neither the name of John MacFarlane nor the names of other
% contributors may be used to endorse or promote products derived
% from this software without specific prior written permission.
%
% THIS SOFTWARE IS PROVIDED BY THE COPYRIGHT HOLDERS AND CONTRIBUTORS
% "AS IS" AND ANY EXPRESS OR IMPLIED WARRANTIES, INCLUDING, BUT NOT
% LIMITED TO, THE IMPLIED WARRANTIES OF MERCHANTABILITY AND FITNESS
% FOR A PARTICULAR PURPOSE ARE DISCLAIMED. IN NO EVENT SHALL THE
% COPYRIGHT OWNER OR CONTRIBUTORS BE LIABLE FOR ANY DIRECT, INDIRECT,
% INCIDENTAL, SPECIAL, EXEMPLARY, OR CONSEQUENTIAL DAMAGES (INCLUDING,
% BUT NOT LIMITED TO, PROCUREMENT OF SUBSTITUTE GOODS OR SERVICES;
% LOSS OF USE, DATA, OR PROFITS; OR BUSINESS INTERRUPTION) HOWEVER
% CAUSED AND ON ANY THEORY OF LIABILITY, WHETHER IN CONTRACT, STRICT
% LIABILITY, OR TORT (INCLUDING NEGLIGENCE OR OTHERWISE) ARISING IN
% ANY WAY OUT OF THE USE OF THIS SOFTWARE, EVEN IF ADVISED OF THE
% POSSIBILITY OF SUCH DAMAGE.
%%

%%
% This is the Eisvogel pandoc LaTeX template.
%
% For usage information and examples visit the official GitHub page:
% https://github.com/Wandmalfarbe/pandoc-latex-template
%%
% Options for packages loaded elsewhere
\PassOptionsToPackage{unicode}{hyperref}
\PassOptionsToPackage{hyphens}{url}
\PassOptionsToPackage{dvipsnames,svgnames,x11names,table}{xcolor}
\documentclass[
  english,
  paper=a4,
  oneside  ,captions=tableheading
]{scrartcl}
\usepackage{xcolor}
\usepackage[margin=2.5cm,includehead=true,includefoot=true,centering,]{geometry}
\usepackage{amsmath,amssymb}

% add backlinks to footnote references, cf. https://tex.stackexchange.com/questions/302266/make-footnote-clickable-both-ways
\usepackage{footnotebackref}
\setcounter{secnumdepth}{5}
\usepackage{iftex}
\ifPDFTeX
  \usepackage[T1]{fontenc}
  \usepackage[utf8]{inputenc}
  \usepackage{textcomp} % provide euro and other symbols
\else % if luatex or xetex
  \usepackage{unicode-math} % this also loads fontspec
  \defaultfontfeatures{Scale=MatchLowercase}
  \defaultfontfeatures[\rmfamily]{Ligatures=TeX,Scale=1}
\fi
\usepackage{lmodern}
\ifPDFTeX\else
  % xetex/luatex font selection
\fi
% Use upquote if available, for straight quotes in verbatim environments
\IfFileExists{upquote.sty}{\usepackage{upquote}}{}
\IfFileExists{microtype.sty}{% use microtype if available
  \usepackage[]{microtype}
  \UseMicrotypeSet[protrusion]{basicmath} % disable protrusion for tt fonts
}{}

% Use setspace anyway because we change the default line spacing.
% The spacing is changed early to affect the titlepage and the TOC.
\usepackage{setspace}
\setstretch{1.2}
\makeatletter
\@ifundefined{KOMAClassName}{% if non-KOMA class
  \IfFileExists{parskip.sty}{%
    \usepackage{parskip}
  }{% else
    \setlength{\parindent}{0pt}
    \setlength{\parskip}{6pt plus 2pt minus 1pt}}
}{% if KOMA class
  \KOMAoptions{parskip=half}}
\makeatother
\usepackage{color}
\usepackage{fancyvrb}
\newcommand{\VerbBar}{|}
\newcommand{\VERB}{\Verb[commandchars=\\\{\}]}
\DefineVerbatimEnvironment{Highlighting}{Verbatim}{commandchars=\\\{\}}
% Add ',fontsize=\small' for more characters per line
\newenvironment{Shaded}{}{}
\newcommand{\AlertTok}[1]{\textcolor[rgb]{1.00,0.00,0.00}{\textbf{#1}}}
\newcommand{\AnnotationTok}[1]{\textcolor[rgb]{0.38,0.63,0.69}{\textbf{\textit{#1}}}}
\newcommand{\AttributeTok}[1]{\textcolor[rgb]{0.49,0.56,0.16}{#1}}
\newcommand{\BaseNTok}[1]{\textcolor[rgb]{0.25,0.63,0.44}{#1}}
\newcommand{\BuiltInTok}[1]{\textcolor[rgb]{0.00,0.50,0.00}{#1}}
\newcommand{\CharTok}[1]{\textcolor[rgb]{0.25,0.44,0.63}{#1}}
\newcommand{\CommentTok}[1]{\textcolor[rgb]{0.38,0.63,0.69}{\textit{#1}}}
\newcommand{\CommentVarTok}[1]{\textcolor[rgb]{0.38,0.63,0.69}{\textbf{\textit{#1}}}}
\newcommand{\ConstantTok}[1]{\textcolor[rgb]{0.53,0.00,0.00}{#1}}
\newcommand{\ControlFlowTok}[1]{\textcolor[rgb]{0.00,0.44,0.13}{\textbf{#1}}}
\newcommand{\DataTypeTok}[1]{\textcolor[rgb]{0.56,0.13,0.00}{#1}}
\newcommand{\DecValTok}[1]{\textcolor[rgb]{0.25,0.63,0.44}{#1}}
\newcommand{\DocumentationTok}[1]{\textcolor[rgb]{0.73,0.13,0.13}{\textit{#1}}}
\newcommand{\ErrorTok}[1]{\textcolor[rgb]{1.00,0.00,0.00}{\textbf{#1}}}
\newcommand{\ExtensionTok}[1]{#1}
\newcommand{\FloatTok}[1]{\textcolor[rgb]{0.25,0.63,0.44}{#1}}
\newcommand{\FunctionTok}[1]{\textcolor[rgb]{0.02,0.16,0.49}{#1}}
\newcommand{\ImportTok}[1]{\textcolor[rgb]{0.00,0.50,0.00}{\textbf{#1}}}
\newcommand{\InformationTok}[1]{\textcolor[rgb]{0.38,0.63,0.69}{\textbf{\textit{#1}}}}
\newcommand{\KeywordTok}[1]{\textcolor[rgb]{0.00,0.44,0.13}{\textbf{#1}}}
\newcommand{\NormalTok}[1]{#1}
\newcommand{\OperatorTok}[1]{\textcolor[rgb]{0.40,0.40,0.40}{#1}}
\newcommand{\OtherTok}[1]{\textcolor[rgb]{0.00,0.44,0.13}{#1}}
\newcommand{\PreprocessorTok}[1]{\textcolor[rgb]{0.74,0.48,0.00}{#1}}
\newcommand{\RegionMarkerTok}[1]{#1}
\newcommand{\SpecialCharTok}[1]{\textcolor[rgb]{0.25,0.44,0.63}{#1}}
\newcommand{\SpecialStringTok}[1]{\textcolor[rgb]{0.73,0.40,0.53}{#1}}
\newcommand{\StringTok}[1]{\textcolor[rgb]{0.25,0.44,0.63}{#1}}
\newcommand{\VariableTok}[1]{\textcolor[rgb]{0.10,0.09,0.49}{#1}}
\newcommand{\VerbatimStringTok}[1]{\textcolor[rgb]{0.25,0.44,0.63}{#1}}
\newcommand{\WarningTok}[1]{\textcolor[rgb]{0.38,0.63,0.69}{\textbf{\textit{#1}}}}

% Workaround/bugfix from jannick0.
% See https://github.com/jgm/pandoc/issues/4302#issuecomment-360669013)
% or https://github.com/Wandmalfarbe/pandoc-latex-template/issues/2
%
% Redefine the verbatim environment 'Highlighting' to break long lines (with
% the help of fvextra). Redefinition is necessary because it is unlikely that
% pandoc includes fvextra in the default template.
\usepackage{fvextra}
\DefineVerbatimEnvironment{Highlighting}{Verbatim}{breaklines,fontsize=\small,commandchars=\\\{\}}

\usepackage{longtable,booktabs,array}
\newcounter{none} % for unnumbered tables
\usepackage{calc} % for calculating minipage widths
% Correct order of tables after \paragraph or \subparagraph
\usepackage{etoolbox}
\makeatletter
\patchcmd\longtable{\par}{\if@noskipsec\mbox{}\fi\par}{}{}
\makeatother
% Allow footnotes in longtable head/foot
\IfFileExists{footnotehyper.sty}{\usepackage{footnotehyper}}{\usepackage{footnote}}
\makesavenoteenv{longtable}
\ifLuaTeX
\usepackage[bidi=basic,shorthands=off]{babel}
\else
\usepackage[bidi=default,shorthands=off]{babel}
\fi
\ifLuaTeX
  \usepackage{selnolig} % disable illegal ligatures
\fi
\setlength{\emergencystretch}{3em} % prevent overfull lines
\providecommand{\tightlist}{%
  \setlength{\itemsep}{0pt}\setlength{\parskip}{0pt}}
\usepackage{fontawesome5}
\usepackage{hyperref}
\hypersetup{
  colorlinks=true,
  linkcolor=blue,
  filecolor=blue,
  urlcolor=blue,
  citecolor=blue,
  pdfborder={0 0 0}
}
\newcommand{\usefullink}[2]{%
  \par\vspace{0.5ex}%
  \href{#1}{\textcolor{blue}{\underline{\faLink\quad #2}}}%
  \par\vspace{0.5ex}%
}
\usepackage{bookmark}
\IfFileExists{xurl.sty}{\usepackage{xurl}}{} % add URL line breaks if available
\urlstyle{same}
\definecolor{default-linkcolor}{HTML}{A50000}
\definecolor{default-filecolor}{HTML}{A50000}
\definecolor{default-citecolor}{HTML}{4077C0}
\definecolor{default-urlcolor}{HTML}{4077C0}

\hypersetup{
  pdflang={en},
  hidelinks,
  breaklinks=true,
  pdfcreator={LaTeX via pandoc with the Eisvogel template}}

\author{}
\date{}


%
% for the background color of the title page
%

%
% break urls
%
\PassOptionsToPackage{hyphens}{url}

%
% When using babel or polyglossia with biblatex, loading csquotes is recommended
% to ensure that quoted texts are typeset according to the rules of your main language.
%
\usepackage{csquotes}

%
% captions
%
\definecolor{caption-color}{HTML}{777777}
\usepackage[font={stretch=1.2}, textfont={color=caption-color}, position=top, skip=4mm, labelfont=bf, singlelinecheck=false, justification=raggedright]{caption}
\setcapindent{0em}

%
% blockquote
%
\definecolor{blockquote-border}{RGB}{221,221,221}
\definecolor{blockquote-text}{RGB}{119,119,119}
\usepackage{mdframed}
\newmdenv[rightline=false,bottomline=false,topline=false,linewidth=3pt,linecolor=blockquote-border,skipabove=\parskip]{customblockquote}
\renewenvironment{quote}{\begin{customblockquote}\list{}{\rightmargin=0em\leftmargin=0em}%
\item\relax\color{blockquote-text}\ignorespaces}{\unskip\unskip\endlist\end{customblockquote}}

%
% Source Sans Pro as the default font family
% Source Code Pro for monospace text
%
% 'default' option sets the default
% font family to Source Sans Pro, not \sfdefault.
%
% Note that the font has been officially renamed to `Source Sans 3`, and
% the version provided by the `sourcesanspro` package is slightly outdated.
% You can install the newer version locally and use it, for example, with
% `mainfont: "Source Sans 3"` in the YAML metadata (requires XeTeX or LuaTeX).
%
\ifnum 0\ifxetex 1\fi\ifluatex 1\fi=0 % if pdftex
    \usepackage[default]{sourcesanspro}
  \usepackage{sourcecodepro}
  \else % if not pdftex
    \usepackage[default]{sourcesanspro}
  \usepackage{sourcecodepro}

  % XeLaTeX specific adjustments for straight quotes: https://tex.stackexchange.com/a/354887
  % This issue is already fixed (see https://github.com/silkeh/latex-sourcecodepro/pull/5) but the
  % fix is still unreleased.
  % TODO: Remove this workaround when the new version of sourcecodepro is released on CTAN.
  \ifxetex
    \makeatletter
    \defaultfontfeatures[\ttfamily]
      { Numbers   = \sourcecodepro@figurestyle,
        Scale     = \SourceCodePro@scale,
        Extension = .otf }
    \setmonofont
      [ UprightFont    = *-\sourcecodepro@regstyle,
        ItalicFont     = *-\sourcecodepro@regstyle It,
        BoldFont       = *-\sourcecodepro@boldstyle,
        BoldItalicFont = *-\sourcecodepro@boldstyle It ]
      {SourceCodePro}
    \makeatother
  \fi
  \fi

%
% heading color
%
\definecolor{heading-color}{RGB}{40,40,40}
% By default, the KOMA-Script classes will typeset sectioning headings in
% sans-serif. Use the normal body font for headings.
\addtokomafont{disposition}{\normalfont\color{heading-color}\bfseries}

%
% variables for title, author and date
%
\usepackage{titling}
\title{}
\author{}
\date{}

%
% tables
%

\definecolor{table-row-color}{HTML}{F5F5F5}
\definecolor{table-rule-color}{HTML}{999999}

%\arrayrulecolor{black!40}
\arrayrulecolor{table-rule-color}     % color of \toprule, \midrule, \bottomrule
\setlength\heavyrulewidth{0.3ex}      % thickness of \toprule, \bottomrule
\renewcommand{\arraystretch}{1.3}     % spacing (padding)


%
% remove paragraph indentation
%
\setlength{\parindent}{0pt}
\setlength{\parskip}{6pt plus 2pt minus 1pt}
\setlength{\emergencystretch}{3em}  % prevent overfull lines

%
%
% Listings
%
%


%
% header and footer
%
\usepackage[headsepline,footsepline]{scrlayer-scrpage}

\newpairofpagestyles{eisvogel-header-footer}{
  \clearpairofpagestyles
  \ihead*{}
  \chead*{}
  \ohead*{}
  \ifoot*{}
  \cfoot*{}
  \ofoot*{\thepage}
  \addtokomafont{pageheadfoot}{\upshape}
}
\pagestyle{eisvogel-header-footer}



%
% Define watermark
%

\begin{document}




\section{Linux Course}\label{linux-course}

\subsection{For AI \& Robotics Applications (Raspberry Pi +
ROS)}\label{for-ai-robotics-applications-raspberry-pi-ros}

This course is a practical Linux foundation for AI and robotics
learners. It is designed for development on Linux machines and VMs, with
Raspberry Pi concepts integrated across sessions.

\textbf{GitHub repository:}
https://github.com/KhaledMahfouz5/linux-course

\textbf{Course structure:} 3 chapters, 11 sessions\\
\textbf{Session duration:} 60-90 minutes each\\
\textbf{Hardware note:} Physical Raspberry Pi is optional

\begin{center}\rule{0.5\linewidth}{0.5pt}\end{center}

\subsection{Requirements}\label{requirements}

\begin{itemize}
\tightlist
\item
  Basic programming knowledge (variables, conditions, loops, functions)
\item
  Internet connection for searching the web and installing packages
  (most of the time you do not need a huge bandwidth)
\item
  30-50 GB free storage
\item
  A note-taking method during class
\end{itemize}

\begin{center}\rule{0.5\linewidth}{0.5pt}\end{center}

\subsection{Chapter 1: Introduction, Installation, and Embedded Linux
Basics}\label{chapter-1-introduction-installation-and-embedded-linux-basics}

\subsubsection{Session 1: Linux for Robotics \&
AI}\label{session-1-linux-for-robotics-ai}

\begin{quote}
\textbf{Before class:} Watch the linked videos for context. Come
curious.
\end{quote}

\paragraph{Learning outcomes}\label{learning-outcomes}

By the end of this session, students should be able to: - Explain why
Linux dominates robotics and AI workflows - Distinguish between the
kernel and a full operating system - Choose suitable Linux distributions
for robotics learning and deployment - Understand the Unix philosophy
and free software movement - Evaluate Linux strengths/limitations for
embedded and AI use-cases - Understand the Raspberry Pi ecosystem at a
conceptual level

\begin{center}\rule{0.5\linewidth}{0.5pt}\end{center}

\paragraph{1) Why Linux Matters Today}\label{why-linux-matters-today}

If you look at the modern world---servers, cloud platforms,
supercomputers, phones, routers, IoT, cybersecurity labs, programming
environments---Linux is \emph{everywhere}. It's the invisible engine
running most of the technologies you interact with daily.

A few facts: - \textbf{100\%} of the world's top 500 supercomputers run
Linux - \textbf{Most web servers} use Linux (Apache, Nginx) -
\textbf{Android}, the most widely-used mobile OS, is built on the Linux
kernel - \textbf{Cybersecurity distributions} (Kali, Parrot) are
Linux-based - Developers choose Linux for transparency, speed,
customizability, and control

So even if you've never installed Linux on your laptop, you've
definitely used it indirectly.

For robotics and AI, this matters because the same OS family runs: -
Development laptops - Edge devices (Pi/Jetson-like boards) - CI/CD
pipelines - Data and model-serving infrastructure

Linux gives one consistent operational model from prototype to
deployment.

\begin{center}\rule{0.5\linewidth}{0.5pt}\end{center}

\paragraph{2) Understanding Linux: What It Really
Is}\label{understanding-linux-what-it-really-is}

One of the biggest misconceptions is thinking Linux is ``that thing that
replaces Windows.'' Linux is not a program. It's not a GUI. It's not
even an operating system by itself.

\textbf{Linux = the kernel}

The kernel is the core program that: - Talks to hardware - Manages
memory - Schedules processes - Handles permissions - Controls I/O

Everything else---the desktop, apps, package manager, and
utilities---comes from the broader free-software ecosystem called
\textbf{GNU}.

Together:\\
\textbf{GNU tools + Linux kernel = GNU/Linux (a complete OS)}

This separation is important because Linux teaches you something deeper
than using an OS---it teaches you \emph{how systems work}.

\begin{center}\rule{0.5\linewidth}{0.5pt}\end{center}

\paragraph{3) Kernel vs OS (embedded
perspective)}\label{kernel-vs-os-embedded-perspective}

\textbf{Kernel} = low-level core that manages CPU, memory, processes,
devices, and system calls.\\
\textbf{Operating system} = kernel + userland tools + services + package
system + optional desktop.

Embedded systems usually optimize the userland while relying on the same
Linux kernel principles: - Hardware abstraction through drivers -
Process isolation - Permission model - Filesystem-based device
interfaces

\begin{center}\rule{0.5\linewidth}{0.5pt}\end{center}

\paragraph{4) The Unix Philosophy}\label{the-unix-philosophy}

Linux inherits its design principles from Unix. This philosophy will
shape how you think during the whole course:

\begin{enumerate}
\def\labelenumi{\arabic{enumi}.}
\tightlist
\item
  \textbf{Write programs that do one thing, and do it well}
\item
  \textbf{Write programs to work together}
\item
  \textbf{Write programs to handle text, because text is a universal
  interface}
\end{enumerate}

This leads to: - Tools that are small and powerful - Commands that do
simple tasks - The magic happens when you combine them using pipes
(\texttt{\textbar{}})

It's why a Linux power-user is unstoppable---every small tool becomes a
Lego block, and you can build anything from those blocks.

\begin{center}\rule{0.5\linewidth}{0.5pt}\end{center}

\paragraph{5) Linux History (The Simple
Version)}\label{linux-history-the-simple-version}

Let's keep it practical.

\begin{itemize}
\tightlist
\item
  \textbf{1969} -- Unix born. AT\&T created Unix. Clean, elegant,
  powerful.
\item
  \textbf{1983} -- GNU Project launched. Richard Stallman began building
  a free Unix-like system. GNU had everything---compilers, libraries,
  tools---\emph{except a kernel}.
\item
  \textbf{1991} -- Linus Torvalds writes the Linux Kernel. As a personal
  hobby project. He released it under a free license, allowing anyone to
  study, modify, and contribute.
\item
  \textbf{1992} -- GNU + Linux merged. This created the first fully free
  operating system.
\item
  \textbf{1990s--2000s} -- Linux takes over servers. Why? Stability,
  control, security, no licensing fees.
\item
  \textbf{2010s--2020s} -- Linux becomes everywhere. Cloud, containers,
  mobile, embedded systems, DevOps\ldots{} The world quietly
  standardized on Linux.
\end{itemize}

That's why learning Linux today is a \textbf{career skill}, not just a
personal preference.

\begin{center}\rule{0.5\linewidth}{0.5pt}\end{center}

\paragraph{6) The Free Software Philosophy (Why Linux
Exists)}\label{the-free-software-philosophy-why-linux-exists}

Linux exists because of the free software movement. Free software is
about \emph{freedom}, not price.

\textbf{The Four Essential Freedoms:} 1. \textbf{Freedom 0:} Run the
program for any purpose 2. \textbf{Freedom 1:} Study the source code and
modify it 3. \textbf{Freedom 2:} Distribute copies 4. \textbf{Freedom
3:} Distribute modified versions

If a program denies any of these freedoms, it becomes
\textbf{proprietary}, which means: - You can't see how it works - You
can't fix it - You depend on the vendor - You lose control over your own
computer

Linux was built to solve this. During this course, you'll see the impact
of this philosophy: - Software you can fully audit - Tools you can
modify - Community-built contributions - Transparency at every layer

This is why cybersecurity professionals prefer Linux .. you can see
\emph{everything} happening in the system.

\begin{center}\rule{0.5\linewidth}{0.5pt}\end{center}

\paragraph{7) Why Robotics Engineers Choose
Linux}\label{why-robotics-engineers-choose-linux}

\begin{quote}
→ \textbf{See Session 6 §4-§5} for deep dive into permissions, root
access, and safe privilege usage.
\end{quote}

\begin{itemize}
\tightlist
\item
  Strong CLI tooling for automation and debugging
\item
  Better compatibility with robotics stacks (ROS ecosystem)
\item
  Native package and build tooling for C/C++/Python workflows
\item
  Reliable remote management over SSH
\item
  Easy scripting for repetitive hardware/software tasks
\item
  Everything is a file---devices, sockets, processes, IPC\ldots{} all
  represented as files. Simplifies development massively
\item
  Powerful terminal environment---you can automate anything
\item
  Package managers---install complete development stacks in seconds
\item
  DevOps and cloud-native tools---Docker, Kubernetes, CI pipelines---all
  built on Linux
\item
  High performance and stability---perfect for servers and embedded
  systems
\item
  Customizability---modify your environment to match exactly how you
  think and work
\end{itemize}

\begin{center}\rule{0.5\linewidth}{0.5pt}\end{center}

\paragraph{8) Linux Distributions for
Robotics}\label{linux-distributions-for-robotics}

If Linux is the engine, a \textbf{distro} is the full car.

\textbf{Ubuntu:} most common ROS learning path, huge documentation, easy
onboarding\\
\textbf{Debian:} very stable and lightweight, strong base for long-term
systems\\
\textbf{Raspberry Pi OS (Raspbian):} Pi-focused integration and
educational accessibility

Selection criteria: - Stability and long-term support - Package
availability - Hardware support - Team familiarity

\textbf{Beginner-friendly:} Ubuntu, Linux Mint, MX Linux\\
\textbf{Power-user:} Arch Linux, Gentoo, Void Linux\\
\textbf{Enterprise:} Red Hat Enterprise Linux, SUSE Linux Enterprise\\
\textbf{Security:} Kali Linux, Parrot OS

\begin{quote}
\textbf{Recommendation:} For this course, use desktop Ubuntu or Debian
for easiest package support and lab consistency. Avoid Arch-based
systems as your first Linux unless you're ready for deeper learning.
\end{quote}

\begin{center}\rule{0.5\linewidth}{0.5pt}\end{center}

\paragraph{9) Pros and Cons of Linux in
Robotics/AI}\label{pros-and-cons-of-linux-in-roboticsai}

\textbf{Pros:} - High control and transparency - Excellent automation
capabilities - Strong networking and remote tooling - Efficient on
modest hardware - Powerful CLI - Free and open source - Fast,
lightweight - Secure and stable - Full control over your system - Great
for programming and networking - Massive community support

\textbf{Cons:} - Higher learning curve for beginners - Some
device/vendor tooling may be Windows-first - Some hardware drivers might
be missing - Gaming is improving but still imperfect - Requires
operational discipline (permissions, services, system config) - Requires
learning mindset, not point-and-click

But as a developer, engineer, or cybersecurity student, the \emph{pros}
outweigh the cons by a huge margin.

\begin{center}\rule{0.5\linewidth}{0.5pt}\end{center}

\paragraph{10) Raspberry Pi Ecosystem Overview
(Conceptual)}\label{raspberry-pi-ecosystem-overview-conceptual}

Key idea: Raspberry Pi is a compact Linux computer often used for
sensing, control, and edge inference.

Ecosystem concepts: - GPIO for digital I/O - Camera and CSI-based vision
capture - I2C/SPI/UART buses for sensors/actuators - SD-card based
storage and boot - SSH-first/headless management patterns

\begin{center}\rule{0.5\linewidth}{0.5pt}\end{center}

\paragraph{11) How to Think Like a Linux
User}\label{how-to-think-like-a-linux-user}

To succeed in this course:

\textbf{Adopt these habits:} - \textbf{Be curious} --- try commands,
explore folders - \textbf{Read error messages} --- Linux errors are
usually helpful - \textbf{Use man pages} --- every command has
documentation - \textbf{Experiment safely} --- break things in virtual
machines - \textbf{Think in small tools} --- pipes, filters, text
processing - \textbf{Love the terminal} --- it's your real power

Linux rewards people who explore.

\begin{center}\rule{0.5\linewidth}{0.5pt}\end{center}

\paragraph{Useful Links --- Session 1}\label{useful-links-session-1}

\textbf{Video Introductions:}

%
  \par\vspace{0.5ex}%
  \href{https://www.youtube.com/watch?v=ShcR4Zfc6Dw}{\textcolor{blue}{\underline{\faLink\quad Linux History \& Distros}}}%
  \par\vspace{0.5ex}%

%
  \par\vspace{0.5ex}%
  \href{https://youtu.be/gojeTqXdBH0?si=kg-iP4KEFjLU-Tto}{\textcolor{blue}{\underline{\faLink\quad Unix Philosophy}}}%
  \par\vspace{0.5ex}%

%
  \par\vspace{0.5ex}%
  \href{https://youtu.be/IvGdY6luTtU?si=nEcir9iqjmYMa1lo}{\textcolor{blue}{\underline{\faLink\quad Kernel vs OS}}}%
  \par\vspace{0.5ex}%

%
  \par\vspace{0.5ex}%
  \href{https://youtu.be/Ag1AKIl_2GM?si=0YDnS2BDSrAAmFX8}{\textcolor{blue}{\underline{\faLink\quad Free Software Philosophy}}}%
  \par\vspace{0.5ex}%

%
  \par\vspace{0.5ex}%
  \href{https://youtu.be/otDOHt\_Jges?si=4ChY6-xpYbOdbKQ6}{\textcolor{blue}{\underline{\faLink\quad Why Programmers Love Linux}}}%
  \par\vspace{0.5ex}%

%
  \par\vspace{0.5ex}%
  \href{https://youtu.be/pYfzZRyRzvs?si=4Q5\_H0gJ3JHBodOc}{\textcolor{blue}{\underline{\faLink\quad Pros and Cons}}}%
  \par\vspace{0.5ex}%

\textbf{Recommended Arabic YouTube Channels:}

%
  \par\vspace{0.5ex}%
  \href{https://www.youtube.com/@anaHr}{\textcolor{blue}{\underline{\faLink\quad anaHr}}}%
  \par\vspace{0.5ex}%

%
  \par\vspace{0.5ex}%
  \href{https://www.youtube.com/@mmbesar}{\textcolor{blue}{\underline{\faLink\quad mohammad besar}}}%
  \par\vspace{0.5ex}%

%
  \par\vspace{0.5ex}%
  \href{https://www.youtube.com/@alwaqad}{\textcolor{blue}{\underline{\faLink\quad Al-Waqqad}}}%
  \par\vspace{0.5ex}%

%
  \par\vspace{0.5ex}%
  \href{https://www.youtube.com/@Sudo\_Start}{\textcolor{blue}{\underline{\faLink\quad sudostart}}}%
  \par\vspace{0.5ex}%

%
  \par\vspace{0.5ex}%
  \href{https://www.youtube.com/@linuxtopia}{\textcolor{blue}{\underline{\faLink\quad linuxtopia}}}%
  \par\vspace{0.5ex}%

%
  \par\vspace{0.5ex}%
  \href{https://www.youtube.com/@abdulmogeeb}{\textcolor{blue}{\underline{\faLink\quad Abdulmojeeb Al-Hameed}}}%
  \par\vspace{0.5ex}%

\textbf{Book:} \emph{The Art of UNIX Programming} --- Eric Raymond

\begin{center}\rule{0.5\linewidth}{0.5pt}\end{center}

\subsubsection{Session 2: Installation \& Embedded
Setup}\label{session-2-installation-embedded-setup}

\begin{quote}
\textbf{Before class:} Install VirtualBox/VMware. Optional: QEMU for Pi
emulation.
\end{quote}

\paragraph{Learning outcomes}\label{learning-outcomes-1}

\begin{itemize}
\tightlist
\item
  Install Linux via VM or dual boot safely
\item
  Understand headless setup basics (SSH/WiFi/static IP)
\item
  Perform post-install essentials for reliable daily use
\item
  Explain the role of Camera/I2C/SPI/UART in Pi-based systems
\end{itemize}

\begin{center}\rule{0.5\linewidth}{0.5pt}\end{center}

\paragraph{1) Choosing Your Linux
Distribution}\label{choosing-your-linux-distribution}

\begin{quote}
→ \textbf{See Session 1 §8} for detailed distro categories (beginner,
power-user, enterprise, security) and selection criteria.
\end{quote}

Before installing Linux, you must pick the right \textbf{distro}
(distribution). Think of distros like different flavors of the same
operating system---they all share the Linux kernel but differ in user
interface, package managers, performance, and philosophy.

\textbf{What You Should Look For:}

{\def\LTcaptype{none} % do not increment counter
\begin{longtable}[]{@{}
  >{\raggedright\arraybackslash}p{(\linewidth - 4\tabcolsep) * \real{0.3333}}
  >{\raggedright\arraybackslash}p{(\linewidth - 4\tabcolsep) * \real{0.3333}}
  >{\raggedright\arraybackslash}p{(\linewidth - 4\tabcolsep) * \real{0.3333}}@{}}
\toprule\noalign{}
\begin{minipage}[b]{\linewidth}\raggedright
Requirement
\end{minipage} & \begin{minipage}[b]{\linewidth}\raggedright
What it means
\end{minipage} & \begin{minipage}[b]{\linewidth}\raggedright
Good Choices
\end{minipage} \\
\midrule\noalign{}
\endhead
\bottomrule\noalign{}
\endlastfoot
\textbf{Ease of use} & Beginner-friendly UI, simple updates & Ubuntu,
Linux Mint, Pop!\_OS \\
\textbf{Stability} & Fewer issues \& long-term support & Ubuntu LTS,
Debian \\
\textbf{Performance} & Works well on older hardware & Xubuntu, Linux
Lite \\
\textbf{Security} & Good by default, strong community & Fedora,
Ubuntu \\
\textbf{Software availability} & Access to apps \& drivers & Ubuntu,
Fedora \\
\end{longtable}
}

\textbf{Quick Tips:} - If your PC is older than \textbf{10 years}, pick
\textbf{Xubuntu}, \textbf{Linux Mint XFCE}, \textbf{LMDE},
\textbf{MX-Linux}, or \textbf{Zorin OS} - If you're into
\textbf{gaming}, choose **Pop!\_OS** or \textbf{Nobara OS} - If you want
a macOS-like UI, choose \textbf{elementaryOS}

\begin{center}\rule{0.5\linewidth}{0.5pt}\end{center}

\paragraph{2) Primary Setup: Full Linux Installation (Dual Boot or
VM)}\label{primary-setup-full-linux-installation-dual-boot-or-vm}

\textbf{Before You Install -- VERY Important:}

\textbf{A. Backup Your Files}\\
The installer will delete the Windows partition if you're switching
fully. Backup all files from Desktop, Documents, Pictures, Music,
Videos. Move them to D:~(or another partition), external HDD, or USB
stick.

\textbf{B. Check Your Boot Mode}\\
Open Start → type \textbf{System Information} → look for \textbf{UEFI}
(modern, secure boot) or \textbf{Legacy BIOS} (older systems). This
affects partitioning and bootloader behavior.

\textbf{C. Create a Bootable USB}\\
Recommended tool: \textbf{Ventoy}---supports multiple ISOs on the same
flash drive, very beginner-friendly.

\textbf{Try Ubuntu Before Installing:}\\
Boot from the USB and choose ``Try Ubuntu'' (Live Mode). Test Wi-Fi,
Bluetooth, sound, display, keyboard, touchpad. If everything works →
safe to install.

\begin{center}\rule{0.5\linewidth}{0.5pt}\end{center}

\textbf{VM Path (Recommended for Beginners):} - Low risk, fast reset,
easy snapshots - Good for command practice and repeatable labs

\textbf{Recommended VirtualBox Settings:} - EFI: Choose What Gives You
The Best Experience . - CPU: 2-4 cpu cores . - Display: 128 MB VRAM -
Network: NAT (Easy For Beginners) - RAM: 4-6 GB allocated

\textbf{Installing Guest Additions:}

\begin{Shaded}
\begin{Highlighting}[]
\FunctionTok{sudo}\NormalTok{ apt install build{-}essential dkms linux{-}headers{-}}\VariableTok{$(}\FunctionTok{uname} \AttributeTok{{-}r}\VariableTok{)}
\FunctionTok{sudo}\NormalTok{ sh /media/username/VBox}\PreprocessorTok{*}\NormalTok{/VBoxLinuxAdditions.run}
\end{Highlighting}
\end{Shaded}

This enables auto-resize display, clipboard sharing, drag \& drop, and
better performance.

\textbf{Dual-Boot Path (Optional):} - Better native performance -
Requires careful partitioning and backup discipline - Choose ``Something
Else (Manual Partitioning)'' .. DO NOT choose ``Erase disk'' or
``Install alongside Windows''

\textbf{Recommended Partition Layout:}

{\def\LTcaptype{none} % do not increment counter
\begin{longtable}[]{@{}
  >{\raggedright\arraybackslash}p{(\linewidth - 8\tabcolsep) * \real{0.2000}}
  >{\raggedright\arraybackslash}p{(\linewidth - 8\tabcolsep) * \real{0.2000}}
  >{\raggedright\arraybackslash}p{(\linewidth - 8\tabcolsep) * \real{0.2000}}
  >{\raggedright\arraybackslash}p{(\linewidth - 8\tabcolsep) * \real{0.2000}}
  >{\raggedright\arraybackslash}p{(\linewidth - 8\tabcolsep) * \real{0.2000}}@{}}
\toprule\noalign{}
\begin{minipage}[b]{\linewidth}\raggedright
Partition
\end{minipage} & \begin{minipage}[b]{\linewidth}\raggedright
Size
\end{minipage} & \begin{minipage}[b]{\linewidth}\raggedright
Filesystem
\end{minipage} & \begin{minipage}[b]{\linewidth}\raggedright
Mountpoint
\end{minipage} & \begin{minipage}[b]{\linewidth}\raggedright
Notes
\end{minipage} \\
\midrule\noalign{}
\endhead
\bottomrule\noalign{}
\endlastfoot
\textbf{EFI/System} & 512--1024 MB & FAT32 & \texttt{/boot/efi} & Needed
only for UEFI \\
\textbf{ROOT} & 40--100 GB+ & EXT4 & \texttt{/} & Main system \\
\textbf{SWAP} & Optional (2--4 GB) & swap & swap & Only if RAM
\textless{} 8GB \\
\textbf{Data} & Keep & NTFS & \texttt{/mnt/data1} & Do \emph{not}
format \\
\end{longtable}
}

\begin{quote}
\textbf{Beginner's Tip:} Don't create a separate \texttt{/home}
partition. You can add it later when you're more experienced.
\end{quote}

\begin{center}\rule{0.5\linewidth}{0.5pt}\end{center}

\paragraph{3) Headless Setup Concepts}\label{headless-setup-concepts}

Headless means ``manage without monitor/keyboard on target device.''

Core concepts: - SSH service enabled on boot - Network reachability
(WiFi/Ethernet) - Static IP plan for stable robot addressing

\begin{center}\rule{0.5\linewidth}{0.5pt}\end{center}

\paragraph{4) Post-Install Essentials}\label{post-install-essentials}

\textbf{Update System:}

\begin{Shaded}
\begin{Highlighting}[]
\FunctionTok{sudo}\NormalTok{ apt update }\KeywordTok{\&\&} \FunctionTok{sudo}\NormalTok{ apt upgrade }\AttributeTok{{-}y}
\end{Highlighting}
\end{Shaded}

\textbf{Install Ubuntu Restricted Extras (codecs for MP3, MP4, MOV,
Fonts):}

\begin{Shaded}
\begin{Highlighting}[]
\FunctionTok{sudo}\NormalTok{ apt install ubuntu{-}restricted{-}extras}
\end{Highlighting}
\end{Shaded}

\textbf{Add Flatpak + Flathub:}

\begin{Shaded}
\begin{Highlighting}[]
\FunctionTok{sudo}\NormalTok{ apt install flatpak}
\ExtensionTok{flatpak}\NormalTok{ remote{-}add }\AttributeTok{{-}{-}if{-}not{-}exists}\NormalTok{ flathub https://flathub.org/repo/flathub.flatpakrepo}
\FunctionTok{sudo}\NormalTok{ apt install gnome{-}software gnome{-}software{-}plugin{-}flatpak}
\end{Highlighting}
\end{Shaded}

\textbf{AppImage Support:}

\begin{Shaded}
\begin{Highlighting}[]
\FunctionTok{sudo}\NormalTok{ apt install libfuse2}
\end{Highlighting}
\end{Shaded}

\textbf{GNOME Extensions \& Tweaks:}

\begin{Shaded}
\begin{Highlighting}[]
\FunctionTok{sudo}\NormalTok{ apt install chrome{-}gnome{-}shell gnome{-}shell{-}extension{-}manager gnome{-}tweaks}
\end{Highlighting}
\end{Shaded}

\textbf{Minimize on Click:}

\begin{Shaded}
\begin{Highlighting}[]
\ExtensionTok{gsettings}\NormalTok{ set org.gnome.shell.extensions.dash{-}to{-}dock click{-}action }\StringTok{\textquotesingle{}focus{-}minimize{-}or{-}appspread\textquotesingle{}}
\end{Highlighting}
\end{Shaded}

\textbf{Enable New Document Menu:}

\begin{Shaded}
\begin{Highlighting}[]
\FunctionTok{touch}\NormalTok{ \textasciitilde{}/Templates/new}
\end{Highlighting}
\end{Shaded}

\textbf{System Backups:} - \textbf{Timeshift} (System Snapshots):
\texttt{sudo\ apt\ install\ timeshift} - \textbf{Déjà Dup} (User files
backup): Built into Ubuntu

\textbf{Essential Security:} - Install drivers and core tooling - Enable
SSH server - Configure firewall baseline - Create backup/snapshot
routine

\textbf{Final System Update:}

\begin{Shaded}
\begin{Highlighting}[]
\FunctionTok{sudo}\NormalTok{ apt update }\KeywordTok{\&\&} \FunctionTok{sudo}\NormalTok{ apt upgrade }\AttributeTok{{-}y} \KeywordTok{\&\&} \ExtensionTok{flatpak}\NormalTok{ update }\AttributeTok{{-}y} \KeywordTok{\&\&} \FunctionTok{sudo}\NormalTok{ snap refresh}
\end{Highlighting}
\end{Shaded}

\begin{center}\rule{0.5\linewidth}{0.5pt}\end{center}

\paragraph{5) Pi Interfaces
(Theoretical)}\label{pi-interfaces-theoretical}

\begin{quote}
→ \textbf{See Session 1 §10} for the full Raspberry Pi ecosystem
overview (GPIO, Camera, I2C, SPI, UART, SD-card, SSH patterns).
\end{quote}

\begin{itemize}
\tightlist
\item
  \textbf{Camera:} image/video input for CV pipelines
\item
  \textbf{I2C:} low-speed bus for many sensors
\item
  \textbf{SPI:} higher-speed peripheral communication
\item
  \textbf{UART:} serial communication for diagnostics/controllers
\end{itemize}

\begin{center}\rule{0.5\linewidth}{0.5pt}\end{center}

\paragraph{Practice}\label{practice}

Run these commands to verify your SSH and firewall setup:

\begin{Shaded}
\begin{Highlighting}[]
\FunctionTok{sudo}\NormalTok{ apt update }\KeywordTok{\&\&} \FunctionTok{sudo}\NormalTok{ apt upgrade }\AttributeTok{{-}y}
\FunctionTok{sudo}\NormalTok{ apt install }\AttributeTok{{-}y}\NormalTok{ openssh{-}server ufw}
\FunctionTok{sudo}\NormalTok{ systemctl enable }\AttributeTok{{-}{-}now}\NormalTok{ ssh}
\FunctionTok{sudo}\NormalTok{ ufw allow OpenSSH}
\FunctionTok{sudo}\NormalTok{ ufw enable}
\ExtensionTok{ip}\NormalTok{ a}
\end{Highlighting}
\end{Shaded}

\begin{center}\rule{0.5\linewidth}{0.5pt}\end{center}

\paragraph{Useful Links --- Session 2}\label{useful-links-session-2}

\textbf{Video Tutorials:}

%
  \par\vspace{0.5ex}%
  \href{https://youtu.be/8yVlJEzq2eg?si=SWcgh\_Bw51Li3lv6}{\textcolor{blue}{\underline{\faLink\quad Choosing Your Distro}}}%
  \par\vspace{0.5ex}%

%
  \par\vspace{0.5ex}%
  \href{https://www.youtube.com/watch?v=C-a5IamFIuM}{\textcolor{blue}{\underline{\faLink\quad Install Linux Step-by-Step}}}%
  \par\vspace{0.5ex}%

%
  \par\vspace{0.5ex}%
  \href{https://youtu.be/XevGfO\_vQJQ?si=OBqaM8UlUeH1wjdQ}{\textcolor{blue}{\underline{\faLink\quad Dual-Boot with Windows}}}%
  \par\vspace{0.5ex}%

%
  \par\vspace{0.5ex}%
  \href{https://youtu.be/wX75Z-4MEoM?si=raCzHOzYJy-GZi9S}{\textcolor{blue}{\underline{\faLink\quad VirtualBox VM Setup}}}%
  \par\vspace{0.5ex}%

\textbf{Written Guides:}

%
  \par\vspace{0.5ex}%
  \href{https://mmbesar.github.io/Tutorials/Ubuntu-24.04-Post/}{\textcolor{blue}{\underline{\faLink\quad Ubuntu 24.04 Post-Install Guide}}}%
  \par\vspace{0.5ex}%

\begin{center}\rule{0.5\linewidth}{0.5pt}\end{center}

\subsubsection{Session 3: Linux Basics for
Robotics}\label{session-3-linux-basics-for-robotics}

\paragraph{Learning outcomes}\label{learning-outcomes-2}

\begin{itemize}
\tightlist
\item
  Navigate Linux confidently
\item
  Use core command-line tools for daily robotics development
\item
  Understand path usage for ROS workspaces
\item
  Explore device/hardware-facing paths (\texttt{/dev}, \texttt{/sys})
\end{itemize}

\begin{center}\rule{0.5\linewidth}{0.5pt}\end{center}

\paragraph{1) Introduction to the Linux Command
Line}\label{introduction-to-the-linux-command-line}

The Linux command line (often called the \textbf{terminal},
\textbf{shell}, or \textbf{CLI}) is a text-based interface that allows
you to interact with your system quickly and efficiently. While
graphical interfaces (GUIs) are great for daily tasks, the command line
provides:

\begin{itemize}
\tightlist
\item
  Faster automation
\item
  Power-user flexibility
\item
  Access to advanced tools
\item
  A universal interface across all distros
\end{itemize}

Most distributions ship with \textbf{bash} as the default shell, though
others (zsh, fish) also exist.

To open a terminal: - \textbf{Ubuntu / Debian}:
\texttt{Ctrl\ +\ Alt\ +\ T} - \textbf{Fedora}: Activities → Terminal -
\textbf{Any distro}: Search ``Terminal''

\begin{center}\rule{0.5\linewidth}{0.5pt}\end{center}

\paragraph{2) System Information and Hardware
Detection}\label{system-information-and-hardware-detection}

Linux provides several commands to check system details:

\begin{Shaded}
\begin{Highlighting}[]
\FunctionTok{uname} \AttributeTok{{-}a}        \CommentTok{\# All system info}
\FunctionTok{uname} \AttributeTok{{-}r}        \CommentTok{\# Kernel version}
\FunctionTok{uname} \AttributeTok{{-}m}        \CommentTok{\# Architecture}
\ExtensionTok{hostnamectl}     \CommentTok{\# OS name, kernel, hardware}
\end{Highlighting}
\end{Shaded}

\textbf{inxi} (advanced---may need installation):

\begin{Shaded}
\begin{Highlighting}[]
\FunctionTok{sudo}\NormalTok{ apt update}
\FunctionTok{sudo}\NormalTok{ apt install inxi}
\ExtensionTok{inxi} \AttributeTok{{-}F}         \CommentTok{\# CPU, GPU, RAM, kernel, drivers}
\end{Highlighting}
\end{Shaded}

\textbf{lsb\_release} (Debian-based):

\begin{Shaded}
\begin{Highlighting}[]
\ExtensionTok{lsb\_release} \AttributeTok{{-}a}
\end{Highlighting}
\end{Shaded}

\begin{center}\rule{0.5\linewidth}{0.5pt}\end{center}

\paragraph{3) Shell Variables and Environment
Configuration}\label{shell-variables-and-environment-configuration}

The shell stores data (variables) temporarily during your session.

\begin{Shaded}
\begin{Highlighting}[]
\BuiltInTok{echo} \VariableTok{$HOME}      \CommentTok{\# Home directory}
\BuiltInTok{echo} \VariableTok{$USER}      \CommentTok{\# Current user}
\BuiltInTok{echo} \VariableTok{$PATH}      \CommentTok{\# Executable search paths}

\VariableTok{NAME}\OperatorTok{=}\StringTok{"Khaled"}   \CommentTok{\# Create variable}
\BuiltInTok{echo} \VariableTok{$NAME}

\BuiltInTok{export} \VariableTok{EDITOR}\OperatorTok{=}\NormalTok{nano   }\CommentTok{\# Export (available to programs)}

\FunctionTok{printenv}        \CommentTok{\# List all variables}
\end{Highlighting}
\end{Shaded}

Environment variables define runtime behavior for tools and scripts.

\begin{center}\rule{0.5\linewidth}{0.5pt}\end{center}

\paragraph{4) File System Navigation}\label{file-system-navigation}

Every Linux filesystem starts from the \textbf{root (\texttt{/})}
directory.

\begin{Shaded}
\begin{Highlighting}[]
\BuiltInTok{pwd}             \CommentTok{\# Print working directory}
\FunctionTok{ls}              \CommentTok{\# List files}
\FunctionTok{ls} \AttributeTok{{-}l}           \CommentTok{\# Long format}
\FunctionTok{ls} \AttributeTok{{-}la}          \CommentTok{\# Include hidden files}
\end{Highlighting}
\end{Shaded}

\textbf{cd -- Change Directory:}

\begin{Shaded}
\begin{Highlighting}[]
\BuiltInTok{cd}\NormalTok{ /home        }\CommentTok{\# Absolute path}
\BuiltInTok{cd}\NormalTok{ Documents    }\CommentTok{\# Relative path}
\BuiltInTok{cd}\NormalTok{ ..           }\CommentTok{\# Parent directory}
\BuiltInTok{cd}\NormalTok{ \textasciitilde{}            }\CommentTok{\# Home directory}
\BuiltInTok{cd} \AttributeTok{{-}}            \CommentTok{\# Previous directory}
\end{Highlighting}
\end{Shaded}

\textbf{Shortcuts:} - \texttt{.} → current directory - \texttt{..} →
parent directory - \texttt{\textasciitilde{}} → home directory -
\texttt{-} → previous directory

\begin{center}\rule{0.5\linewidth}{0.5pt}\end{center}

\paragraph{5) Essential Commands}\label{essential-commands}

\texttt{ls}, \texttt{cd}, \texttt{pwd}, \texttt{cat}, \texttt{mkdir},
\texttt{cp}, \texttt{mv}, \texttt{rm}, \texttt{echo}, \texttt{uname},
\texttt{inxi} are the baseline daily toolkit.

\begin{Shaded}
\begin{Highlighting}[]
\FunctionTok{mkdir} \AttributeTok{{-}p}\NormalTok{ robotics\_ws/src}
\BuiltInTok{cd}\NormalTok{ robotics\_ws}
\BuiltInTok{pwd}
\BuiltInTok{echo} \StringTok{"session3"} \OperatorTok{\textgreater{}}\NormalTok{ notes.txt}
\FunctionTok{cat}\NormalTok{ notes.txt}
\FunctionTok{cp}\NormalTok{ notes.txt notes.bak}
\FunctionTok{mv}\NormalTok{ notes.bak archive.txt}
\FunctionTok{touch}\NormalTok{ newfile.txt}
\FunctionTok{sleep}\NormalTok{ 5}
\end{Highlighting}
\end{Shaded}

\textbf{rm -- Remove (\faExclamationCircle Dangerous but essential):}

\begin{Shaded}
\begin{Highlighting}[]
\FunctionTok{rm}\NormalTok{ file.txt}
\FunctionTok{rm} \AttributeTok{{-}r}\NormalTok{ folder/}
\FunctionTok{rm} \AttributeTok{{-}rf}\NormalTok{ folder/    }\CommentTok{\# VERY DANGEROUS {-} use with care!}
\end{Highlighting}
\end{Shaded}

\begin{center}\rule{0.5\linewidth}{0.5pt}\end{center}

\paragraph{6) Paths: Relative vs
Absolute}\label{paths-relative-vs-absolute}

Robotics workspaces fail often because of path mistakes.

\textbf{Absolute Paths} start from the root \texttt{/}:

\begin{verbatim}
/home/khaled/projects/code.py
\end{verbatim}

\textbf{Relative Paths} start from your current location:

\begin{verbatim}
../images/photo.png
./run.sh
\end{verbatim}

\begin{itemize}
\tightlist
\item
  \textbf{Absolute:} starts from \texttt{/} (stable in scripts)
\item
  \textbf{Relative:} from current directory (fast interactive use)
\end{itemize}

\begin{center}\rule{0.5\linewidth}{0.5pt}\end{center}

\paragraph{\texorpdfstring{7) Package Managers: \texttt{apt},
\texttt{pip},
\texttt{snap}}{7) Package Managers: apt, pip, snap}}\label{package-managers-apt-pip-snap}

\textbf{apt -- Debian, Ubuntu, Linux Mint:}

\begin{Shaded}
\begin{Highlighting}[]
\FunctionTok{sudo}\NormalTok{ apt update             }\CommentTok{\# Update repositories}
\FunctionTok{sudo}\NormalTok{ apt install gedit      }\CommentTok{\# Install package}
\FunctionTok{sudo}\NormalTok{ apt remove gedit       }\CommentTok{\# Remove package}
\FunctionTok{sudo}\NormalTok{ apt upgrade            }\CommentTok{\# Upgrade system}
\end{Highlighting}
\end{Shaded}

\textbf{dnf -- Fedora, RHEL 8+:}

\begin{Shaded}
\begin{Highlighting}[]
\FunctionTok{sudo}\NormalTok{ dnf install nano}
\FunctionTok{sudo}\NormalTok{ dnf remove nano}
\FunctionTok{sudo}\NormalTok{ dnf update}
\end{Highlighting}
\end{Shaded}

\textbf{yum -- Older RHEL/CentOS:}

\begin{Shaded}
\begin{Highlighting}[]
\FunctionTok{sudo}\NormalTok{ yum install nano}
\FunctionTok{sudo}\NormalTok{ yum update}
\end{Highlighting}
\end{Shaded}

\begin{itemize}
\tightlist
\item
  \texttt{apt} for system packages
\item
  \texttt{pip} for Python packages
\item
  \texttt{snap} for bundled apps/tooling where appropriate
\end{itemize}

\begin{center}\rule{0.5\linewidth}{0.5pt}\end{center}

\paragraph{\texorpdfstring{8) Editors: \texttt{nano}, \texttt{gedit},
\texttt{vim}
Basics}{8) Editors: nano, gedit, vim Basics}}\label{editors-nano-gedit-vim-basics}

\textbf{nano (Terminal Editor):}

\begin{Shaded}
\begin{Highlighting}[]
\FunctionTok{nano}\NormalTok{ file.txt}
\end{Highlighting}
\end{Shaded}

Controls: - Save: \textbf{Ctrl + S} - Exit: \textbf{Ctrl + X} - Search:
\textbf{Ctrl + W}

\textbf{gedit (GUI Editor):}

\begin{Shaded}
\begin{Highlighting}[]
\ExtensionTok{gedit}\NormalTok{ notes.txt }\KeywordTok{\&}    \CommentTok{\# \& lets the terminal continue}
\end{Highlighting}
\end{Shaded}

Students should be comfortable editing files both in terminal and GUI
contexts.

\begin{center}\rule{0.5\linewidth}{0.5pt}\end{center}

\paragraph{9) Help Tools}\label{help-tools}

\begin{quote}
→ \textbf{See Session 1 §11} for the mindset behind using documentation
and self-learning in Linux.
\end{quote}

Linux has multiple built-in help resources.

\begin{Shaded}
\begin{Highlighting}[]
\FunctionTok{man}\NormalTok{ ls              }\CommentTok{\# Manual pages}
\FunctionTok{man}\NormalTok{ tar}
\ExtensionTok{info}\NormalTok{ coreutils      }\CommentTok{\# Info documentation}
\ExtensionTok{tldr}\NormalTok{ ls             }\CommentTok{\# Simplified examples (install first)}
\FunctionTok{ls} \AttributeTok{{-}{-}help}           \CommentTok{\# Quick help}
\FunctionTok{apropos}\NormalTok{ network     }\CommentTok{\# Search commands by description}
\end{Highlighting}
\end{Shaded}

Use \texttt{man}, \texttt{info}, \texttt{tldr}, \texttt{-\/-help}, and
\texttt{apropos} to self-debug and self-learn.

\begin{center}\rule{0.5\linewidth}{0.5pt}\end{center}

\paragraph{\texorpdfstring{10) Robotics Context: \texttt{/dev/} and
\texttt{/sys/}}{10) Robotics Context: /dev/ and /sys/}}\label{robotics-context-dev-and-sys}

In Linux: - Files - Directories - Devices - Hardware - Network
interfaces

→ \textbf{Everything is a file}

\textbf{\texttt{/dev/} (device files):}

\begin{Shaded}
\begin{Highlighting}[]
\ExtensionTok{/dev/sda}          \CommentTok{\# Hard drive}
\ExtensionTok{/dev/null}         \CommentTok{\# Black hole (silence output)}
\ExtensionTok{/dev/tty}          \CommentTok{\# Current terminal}
\ExtensionTok{/dev/random}       \CommentTok{\# Random number generator}
\end{Highlighting}
\end{Shaded}

\textbf{\texttt{/sys/} and \texttt{/proc/} (kernel-exposed
hardware/state metadata):}

\begin{Shaded}
\begin{Highlighting}[]
\ExtensionTok{/proc/cpuinfo}          \CommentTok{\# CPU info}
\ExtensionTok{/proc/uptime}           \CommentTok{\# System uptime}
\ExtensionTok{/proc/net/tcp}          \CommentTok{\# TCP connections}
\ExtensionTok{/sys/class/net/eth0/address}           \CommentTok{\# Network MAC address}
\ExtensionTok{/sys/class/backlight/brightness}       \CommentTok{\# Screen brightness}
\ExtensionTok{/sys/class/power\_supply/BAT0/capacity} \CommentTok{\# Battery level}
\end{Highlighting}
\end{Shaded}

Want hardware info? \texttt{cat\ /proc/cpuinfo}\\
Want to silence output? \texttt{command\ \textgreater{}\ /dev/null}

This design is WHY Linux is scriptable, automatable, and runs the
internet.

\begin{center}\rule{0.5\linewidth}{0.5pt}\end{center}

\paragraph{11) Explore Virtual Hardware Seen by
VM}\label{explore-virtual-hardware-seen-by-vm}

Inspect which virtual disks, network interfaces, and pseudo-devices
appear.

\begin{center}\rule{0.5\linewidth}{0.5pt}\end{center}

\paragraph{\texorpdfstring{12) Archiving: \texttt{tar},
\texttt{zip}}{12) Archiving: tar, zip}}\label{archiving-tar-zip}

Archive projects and logs for sharing, reproducibility, and backups.

\textbf{tar (most common):}

\begin{Shaded}
\begin{Highlighting}[]
\FunctionTok{tar} \AttributeTok{{-}cvf}\NormalTok{ archive.tar folder/           }\CommentTok{\# Create}
\FunctionTok{tar} \AttributeTok{{-}xvf}\NormalTok{ archive.tar                   }\CommentTok{\# Extract}
\FunctionTok{tar} \AttributeTok{{-}czvf}\NormalTok{ archive.tar.gz folder/       }\CommentTok{\# Compressed (gzip)}
\FunctionTok{tar} \AttributeTok{{-}xzvf}\NormalTok{ archive.tar.gz               }\CommentTok{\# Extract compressed}
\end{Highlighting}
\end{Shaded}

\textbf{zip:}

\begin{Shaded}
\begin{Highlighting}[]
\FunctionTok{zip} \AttributeTok{{-}r}\NormalTok{ robotics\_ws.zip robotics\_ws}
\FunctionTok{unzip}\NormalTok{ archive.zip}
\end{Highlighting}
\end{Shaded}

\textbf{7z (7-Zip):}

\begin{Shaded}
\begin{Highlighting}[]
\FunctionTok{sudo}\NormalTok{ apt install p7zip{-}full}
\ExtensionTok{7z}\NormalTok{ a archive.7z folder/}
\ExtensionTok{7z}\NormalTok{ x archive.7z}
\end{Highlighting}
\end{Shaded}

\begin{center}\rule{0.5\linewidth}{0.5pt}\end{center}

\paragraph{Bonus Useful Commands}\label{bonus-useful-commands}

\begin{itemize}
\tightlist
\item
  \texttt{clear} -- clears the terminal
\item
  \texttt{history} -- shows previous commands
\item
  \texttt{head} / \texttt{tail} -- preview files
  (\texttt{tail\ -f\ log.txt} for live logs)
\item
  \texttt{less} -- view files with navigation
\item
  \texttt{grep} -- search text
\item
  \texttt{whoami} -- current user
\item
  \texttt{df\ -h} -- disk usage
\item
  \texttt{du\ -sh\ folder/} -- folder size
\end{itemize}

\begin{center}\rule{0.5\linewidth}{0.5pt}\end{center}

\paragraph{Practice Exercises}\label{practice-exercises}

\begin{enumerate}
\def\labelenumi{\arabic{enumi}.}
\tightlist
\item
  Create a directory structure: \texttt{projects/python/basics}
\item
  Create three files inside \texttt{basics}
\item
  Write your name into \texttt{info.txt} using \texttt{echo}
\item
  Compress the folder using tar and zip
\item
  Install \texttt{tldr} and test 5 commands
\item
  View your system kernel version and CPU info using two methods
\end{enumerate}

\begin{center}\rule{0.5\linewidth}{0.5pt}\end{center}

\subsection{Chapter 2: Programming, Development Tools \& ROS
Foundation}\label{chapter-2-programming-development-tools-ros-foundation}

\subsubsection{Session 4: Text Manipulation \& Log
Analysis}\label{session-4-text-manipulation-log-analysis}

\begin{quote}
Welcome to one of the most \emph{useful and exciting} sessions in your
Linux journey. In this session, we learn how Linux becomes a
\textbf{superpower} when working with text, logs, automation, scripting,
and filtering.
\end{quote}

\paragraph{Learning outcomes}\label{learning-outcomes-3}

\begin{itemize}
\tightlist
\item
  Filter and transform large text/log outputs efficiently
\item
  Extract structured fields from sensor logs
\item
  Locate files quickly in large ROS-style directories
\item
  Chain commands with pipes/redirection
\end{itemize}

\begin{center}\rule{0.5\linewidth}{0.5pt}\end{center}

\paragraph{\texorpdfstring{1) \texttt{grep} --- The Search Engine of the
Terminal}{1) grep --- The Search Engine of the Terminal}}\label{grep-the-search-engine-of-the-terminal}

\texttt{grep} stands for \textbf{Global Regular Expression Print}. In
simple words: \textbf{grep searches for text inside files or command
output.}

\begin{Shaded}
\begin{Highlighting}[]
\FunctionTok{grep} \StringTok{"error"}\NormalTok{ /var/log/syslog                }\CommentTok{\# Find lines with "error"}
\FunctionTok{grep} \AttributeTok{{-}i} \StringTok{"linux"}\NormalTok{ notes.txt                   }\CommentTok{\# Ignore case}
\FunctionTok{grep} \AttributeTok{{-}r} \StringTok{"password"}\NormalTok{ /etc                     }\CommentTok{\# Recursive search}
\FunctionTok{grep} \AttributeTok{{-}n} \StringTok{"TODO"}\NormalTok{ project.py                   }\CommentTok{\# Show line numbers}
\FunctionTok{grep} \AttributeTok{{-}v} \StringTok{"DEBUG"}\NormalTok{ logs.txt                    }\CommentTok{\# Invert match (NOT containing)}
\FunctionTok{grep} \AttributeTok{{-}E} \StringTok{"cat|dog"}\NormalTok{ animals.txt               }\CommentTok{\# Extended regex (OR)}
\FunctionTok{grep} \AttributeTok{{-}i} \StringTok{"error"}\NormalTok{ robot.log}
\FunctionTok{grep} \AttributeTok{{-}nE} \StringTok{"warn|fail|timeout"}\NormalTok{ robot.log}
\end{Highlighting}
\end{Shaded}

\textbf{Practical examples:}

\begin{Shaded}
\begin{Highlighting}[]
\FunctionTok{ls} \KeywordTok{|} \FunctionTok{grep} \StringTok{".pdf"}                            \CommentTok{\# Find PDF files}
\FunctionTok{grep} \StringTok{"failed"} \PreprocessorTok{*}\NormalTok{.log                         }\CommentTok{\# Search in multiple logs}
\FunctionTok{grep} \AttributeTok{{-}E} \StringTok{"[0{-}9]+\textbackslash{}.[0{-}9]+\textbackslash{}.[0{-}9]+\textbackslash{}.[0{-}9]+"}\NormalTok{ nginx.log  }\CommentTok{\# Find IP addresses}
\end{Highlighting}
\end{Shaded}

\begin{center}\rule{0.5\linewidth}{0.5pt}\end{center}

\paragraph{\texorpdfstring{2) \texttt{sed} --- The Stream
Editor}{2) sed --- The Stream Editor}}\label{sed-the-stream-editor}

\texttt{sed} is a \textbf{text transformer}. You can use it to replace
words, delete lines, insert text, extract text, and edit files
\emph{without opening them}.

\begin{Shaded}
\begin{Highlighting}[]
\FunctionTok{sed} \StringTok{\textquotesingle{}s/old/new/\textquotesingle{}}\NormalTok{ file                       }\CommentTok{\# Replace first occurrence per line}
\FunctionTok{sed} \StringTok{\textquotesingle{}s/old/new/g\textquotesingle{}}\NormalTok{ file                      }\CommentTok{\# Replace ALL occurrences}
\FunctionTok{sed} \StringTok{\textquotesingle{}s/linux/Linux/\textquotesingle{}}\NormalTok{ notes.txt}
\FunctionTok{sed} \StringTok{\textquotesingle{}/error/d\textquotesingle{}}\NormalTok{ logs.txt                     }\CommentTok{\# Delete lines with "error"}
\FunctionTok{sed} \AttributeTok{{-}n} \StringTok{\textquotesingle{}/INFO/p\textquotesingle{}}\NormalTok{ logs.txt                  }\CommentTok{\# Print only matching lines}
\FunctionTok{sed} \AttributeTok{{-}i} \StringTok{\textquotesingle{}s/foo/bar/g\textquotesingle{}}\NormalTok{ file.txt              }\CommentTok{\# \textbackslash{}faExclamationCircle Edit IN{-}PLACE (BE CAREFUL!)}
\end{Highlighting}
\end{Shaded}

\textbf{Cool examples:}

\begin{Shaded}
\begin{Highlighting}[]
\FunctionTok{sed} \StringTok{\textquotesingle{}s/ //g\textquotesingle{}}\NormalTok{ text.txt                       }\CommentTok{\# Remove all spaces}
\FunctionTok{sed} \StringTok{\textquotesingle{}s/  */ /g\textquotesingle{}}                             \CommentTok{\# Multiple spaces to one}
\FunctionTok{sed} \StringTok{\textquotesingle{}s/\^{}/[INFO] /\textquotesingle{}}\NormalTok{ file.txt                }\CommentTok{\# Add prefix to each line}
\FunctionTok{sed} \StringTok{\textquotesingle{}s/WARN/WARNING/g\textquotesingle{}}\NormalTok{ robot.log }\OperatorTok{\textgreater{}}\NormalTok{ robot\_clean.log}
\FunctionTok{sed} \AttributeTok{{-}n} \StringTok{\textquotesingle{}/camera/p\textquotesingle{}}\NormalTok{ robot.log}
\end{Highlighting}
\end{Shaded}

\begin{center}\rule{0.5\linewidth}{0.5pt}\end{center}

\paragraph{\texorpdfstring{3) \texttt{awk} --- The Text Processor
King}{3) awk --- The Text Processor King}}\label{awk-the-text-processor-king}

\texttt{awk} is a \textbf{programming language} hidden inside a terminal
command. Perfect for extracting columns, summing numbers, processing CSV
files, and analyzing logs.

\begin{Shaded}
\begin{Highlighting}[]
\FunctionTok{awk} \StringTok{\textquotesingle{}\{print $1\}\textquotesingle{}}\NormalTok{ file.txt                   }\CommentTok{\# Print first column}
\FunctionTok{awk} \StringTok{\textquotesingle{}\{print $2\}\textquotesingle{}}\NormalTok{ data.txt                   }\CommentTok{\# Print second column}
\FunctionTok{awk} \StringTok{\textquotesingle{}$2 \textgreater{} 50\textquotesingle{}}\NormalTok{ data.txt                      }\CommentTok{\# Lines where 2nd field \textgreater{} 50}
\FunctionTok{awk} \StringTok{\textquotesingle{}\{sum += $3\} END \{print sum\}\textquotesingle{}}\NormalTok{ sales.txt }\CommentTok{\# Sum column 3}
\FunctionTok{awk} \AttributeTok{{-}F}\StringTok{\textquotesingle{},\textquotesingle{}} \StringTok{\textquotesingle{}\{print $1, $3\}\textquotesingle{}}\NormalTok{ sensor.csv       }\CommentTok{\# CSV with comma delimiter}
\FunctionTok{awk} \StringTok{\textquotesingle{}$2 \textgreater{} 50 \{print $0\}\textquotesingle{}}\NormalTok{ telemetry.txt}
\end{Highlighting}
\end{Shaded}

\textbf{Example---given line \texttt{Alice\ 23\ Student}:} -
\texttt{\$1} = Alice - \texttt{\$2} = 23 - \texttt{\$3} = Student

\begin{center}\rule{0.5\linewidth}{0.5pt}\end{center}

\paragraph{\texorpdfstring{4) \texttt{find} --- The Ultimate Search
Tool}{4) find --- The Ultimate Search Tool}}\label{find-the-ultimate-search-tool}

\texttt{find} searches your filesystem \textbf{recursively} for files,
directories, sizes, permissions, names, types---everything. Essential in
nested workspaces and mixed code/config trees.

\begin{Shaded}
\begin{Highlighting}[]
\FunctionTok{find}\NormalTok{ . }\AttributeTok{{-}name} \StringTok{"*.txt"}                        \CommentTok{\# Search by name}
\FunctionTok{find}\NormalTok{ /home }\AttributeTok{{-}iname} \StringTok{"linux"}                   \CommentTok{\# Case{-}insensitive}
\FunctionTok{find}\NormalTok{ . }\AttributeTok{{-}type}\NormalTok{ d }\AttributeTok{{-}name} \StringTok{"src"}                  \CommentTok{\# Directories only}
\FunctionTok{find}\NormalTok{ . }\AttributeTok{{-}size}\NormalTok{ +1M                            }\CommentTok{\# By size}
\FunctionTok{find}\NormalTok{ / }\AttributeTok{{-}perm}\NormalTok{ 777                            }\CommentTok{\# By permissions}
\FunctionTok{find}\NormalTok{ . ! }\AttributeTok{{-}name} \StringTok{"*.txt"}                      \CommentTok{\# Exclude pattern}
\FunctionTok{find}\NormalTok{ \textasciitilde{}/robotics\_ws }\AttributeTok{{-}name} \StringTok{"*.launch"}
\FunctionTok{find}\NormalTok{ \textasciitilde{}/robotics\_ws }\AttributeTok{{-}type}\NormalTok{ f }\AttributeTok{{-}name} \StringTok{"*.py"}
\end{Highlighting}
\end{Shaded}

\begin{center}\rule{0.5\linewidth}{0.5pt}\end{center}

\paragraph{5) Pipes and Redirections}\label{pipes-and-redirections}

\textbf{Pipes (\texttt{\textbar{}})} send the output of one command into
another. Think of it as plugging commands together like LEGO pieces.

\begin{Shaded}
\begin{Highlighting}[]
\FunctionTok{ls} \KeywordTok{|} \FunctionTok{grep} \StringTok{".txt"}                            \CommentTok{\# Filter listing}
\FunctionTok{ls} \KeywordTok{|} \FunctionTok{grep} \StringTok{".py"} \KeywordTok{|} \FunctionTok{wc} \AttributeTok{{-}l}                     \CommentTok{\# Count Python files}
\FunctionTok{du} \AttributeTok{{-}ah} \KeywordTok{|} \FunctionTok{sort} \AttributeTok{{-}h} \KeywordTok{|} \FunctionTok{tail}                     \CommentTok{\# Top 10 largest files}
\FunctionTok{cat}\NormalTok{ text.txt }\KeywordTok{|} \FunctionTok{tr} \StringTok{\textquotesingle{} \textquotesingle{}} \StringTok{\textquotesingle{}\textbackslash{}n\textquotesingle{}} \KeywordTok{|} \FunctionTok{sort} \KeywordTok{|} \FunctionTok{uniq}   \CommentTok{\# Unique words}
\end{Highlighting}
\end{Shaded}

\textbf{Redirections:} - \texttt{\textbar{}} passes output between
commands - \texttt{\textgreater{}} overwrite output file -
\texttt{\textgreater{}\textgreater{}} append output file -
\texttt{\textless{}} feed file as input

\begin{Shaded}
\begin{Highlighting}[]
\BuiltInTok{echo} \StringTok{"Hello"} \OperatorTok{\textgreater{}}\NormalTok{ file.txt                     }\CommentTok{\# Write to file}
\BuiltInTok{echo} \StringTok{"World"} \OperatorTok{\textgreater{}\textgreater{}}\NormalTok{ file.txt                    }\CommentTok{\# Append to file}
\FunctionTok{ls} \OperatorTok{\textgreater{}}\NormalTok{ list.txt                               }\CommentTok{\# Redirect output}
\FunctionTok{wc} \AttributeTok{{-}l} \OperatorTok{\textless{}}\NormalTok{ notes.txt                           }\CommentTok{\# Input from file}
\end{Highlighting}
\end{Shaded}

\begin{center}\rule{0.5\linewidth}{0.5pt}\end{center}

\paragraph{6) Practical Robotics
Application}\label{practical-robotics-application}

Filter ROS-like outputs, isolate anomalies, and summarize key fields for
debugging.

\begin{Shaded}
\begin{Highlighting}[]
\FunctionTok{cat}\NormalTok{ robot.log }\KeywordTok{|} \FunctionTok{grep} \StringTok{"camera"} \KeywordTok{|} \FunctionTok{awk} \StringTok{\textquotesingle{}\{print $1,$2,$NF\}\textquotesingle{}}
\FunctionTok{grep} \AttributeTok{{-}i} \StringTok{"imu"}\NormalTok{ topics.txt }\KeywordTok{|} \FunctionTok{sort} \KeywordTok{|} \FunctionTok{uniq} \AttributeTok{{-}c}
\end{Highlighting}
\end{Shaded}

\begin{center}\rule{0.5\linewidth}{0.5pt}\end{center}

\paragraph{Useful Links --- Session 4}\label{useful-links-session-4}

\textbf{Video Tutorials:}

%
  \par\vspace{0.5ex}%
  \href{https://youtu.be/N05sWPgj-44?si=RqN-cKCBSae4bvK7}{\textcolor{blue}{\underline{\faLink\quad grep}}}%
  \par\vspace{0.5ex}%

%
  \par\vspace{0.5ex}%
  \href{https://youtu.be/EACe7aiGczw?si=q8Hyf9d1VgsO-uj2}{\textcolor{blue}{\underline{\faLink\quad sed}}}%
  \par\vspace{0.5ex}%

%
  \par\vspace{0.5ex}%
  \href{https://youtu.be/9YOZmI-zWok?si=Larx5zKGgMMSrsn-}{\textcolor{blue}{\underline{\faLink\quad awk}}}%
  \par\vspace{0.5ex}%

%
  \par\vspace{0.5ex}%
  \href{https://youtu.be/BZ5gsFiIKOQ?si=WEDUXPu6DCi2anNL}{\textcolor{blue}{\underline{\faLink\quad find}}}%
  \par\vspace{0.5ex}%

\begin{center}\rule{0.5\linewidth}{0.5pt}\end{center}

\subsubsection{Session 5: Development Environment for AI \&
Robotics}\label{session-5-development-environment-for-ai-robotics}

\begin{quote}
\textbf{Verify installed:} \texttt{python3}, \texttt{python3-venv},
\texttt{git}, \texttt{cmake}, \texttt{build-essential}
\end{quote}

\paragraph{Learning outcomes}\label{learning-outcomes-4}

\begin{itemize}
\tightlist
\item
  Set up isolated Python development environments
\item
  Install and manage AI/vision packages
\item
  Understand source-build basics
\item
  Use team Git workflows
\item
  Prepare editor + Docker-based development flow
\end{itemize}

\begin{center}\rule{0.5\linewidth}{0.5pt}\end{center}

\paragraph{1) Mindset \& Setup}\label{mindset-setup}

Linux for software dev is about \textbf{three ideas}: 1. Everything is a
tool 2. The terminal is your superpower 3. You \emph{own} your system

\textbf{Update first. Always:}

\begin{Shaded}
\begin{Highlighting}[]
\FunctionTok{sudo}\NormalTok{ apt update }\KeywordTok{\&\&} \FunctionTok{sudo}\NormalTok{ apt upgrade }\AttributeTok{{-}y}
\end{Highlighting}
\end{Shaded}

\textbf{Install the essentials:}

\begin{Shaded}
\begin{Highlighting}[]
\FunctionTok{sudo}\NormalTok{ apt install }\AttributeTok{{-}y}\NormalTok{ build{-}essential curl wget git software{-}properties{-}common}
\end{Highlighting}
\end{Shaded}

\begin{center}\rule{0.5\linewidth}{0.5pt}\end{center}

\paragraph{\texorpdfstring{2) Python Setup (\texttt{venv},
\texttt{pip})}{2) Python Setup (venv, pip)}}\label{python-setup-venv-pip}

Use virtual environments to avoid global dependency conflicts.
\textbf{DO THIS.}

\begin{Shaded}
\begin{Highlighting}[]
\FunctionTok{sudo}\NormalTok{ apt install }\AttributeTok{{-}y}\NormalTok{ python3 python3{-}pip python3{-}venv}
\ExtensionTok{python3} \AttributeTok{{-}m}\NormalTok{ venv .venv}
\BuiltInTok{source}\NormalTok{ .venv/bin/activate}
\ExtensionTok{pip}\NormalTok{ install }\AttributeTok{{-}{-}upgrade}\NormalTok{ pip}
\ExtensionTok{pip}\NormalTok{ install requests bs4}
\end{Highlighting}
\end{Shaded}

\begin{center}\rule{0.5\linewidth}{0.5pt}\end{center}

\paragraph{3) AI/ML Framework
Installation}\label{aiml-framework-installation}

Start with OpenCV and expand to TensorFlow/PyTorch according to hardware
and project scope.

\begin{Shaded}
\begin{Highlighting}[]
\ExtensionTok{pip}\NormalTok{ install opencv{-}python}
\end{Highlighting}
\end{Shaded}

\begin{center}\rule{0.5\linewidth}{0.5pt}\end{center}

\paragraph{4) Compile from Source (Makefiles/Build
Systems)}\label{compile-from-source-makefilesbuild-systems}

Understand compile-link pipeline and why build systems are central in
robotics C/C++ codebases.

\textbf{C/C++ Compilers \& Debugging:}

\begin{Shaded}
\begin{Highlighting}[]
\FunctionTok{sudo}\NormalTok{ apt install }\AttributeTok{{-}y}\NormalTok{ gcc g++ gdb}
\FunctionTok{gcc}\NormalTok{ main.c }\AttributeTok{{-}o}\NormalTok{ main         }\CommentTok{\# Compile}
\FunctionTok{gdb}\NormalTok{ ./main                 }\CommentTok{\# Debug}
\end{Highlighting}
\end{Shaded}

\textbf{Makefile basics:}

\begin{Shaded}
\begin{Highlighting}[]
\DecValTok{all:}
\ErrorTok{    }\NormalTok{gcc main.c {-}o main}
\end{Highlighting}
\end{Shaded}

\begin{Shaded}
\begin{Highlighting}[]
\FunctionTok{make}                       \CommentTok{\# Run Makefile}
\end{Highlighting}
\end{Shaded}

\begin{center}\rule{0.5\linewidth}{0.5pt}\end{center}

\paragraph{5) Git Workflow for Robotics
Teams}\label{git-workflow-for-robotics-teams}

Branch, commit, push, and review with reproducible history and rollback
capability.

\begin{center}\rule{0.5\linewidth}{0.5pt}\end{center}

\paragraph{6) Web Development \& Servers}\label{web-development-servers}

\textbf{Simple HTTP server (classic):}

\begin{Shaded}
\begin{Highlighting}[]
\ExtensionTok{python3} \AttributeTok{{-}m}\NormalTok{ http.server 8000}
\end{Highlighting}
\end{Shaded}

Now anyone on your network can access \texttt{http://YOUR\_IP:8000}

\begin{quote}
\faExclamationCircle Warning: \texttt{http.server} is not recommended
for production. It only implements basic security checks.
\end{quote}

\textbf{Node.js:}

\begin{Shaded}
\begin{Highlighting}[]
\CommentTok{\# Install nvm (Node Version Manager)}
\ExtensionTok{curl} \AttributeTok{{-}o{-}}\NormalTok{ https://raw.githubusercontent.com/nvm{-}sh/nvm/v0.40.3/install.sh }\KeywordTok{|} \FunctionTok{bash}
\ExtensionTok{\textbackslash{}.} \StringTok{"}\VariableTok{$HOME}\StringTok{/.nvm/nvm.sh"}
\ExtensionTok{nvm}\NormalTok{ install 24}

\CommentTok{\# Verify}
\ExtensionTok{node} \AttributeTok{{-}v}
\ExtensionTok{npm} \AttributeTok{{-}v}
\end{Highlighting}
\end{Shaded}

\textbf{PHP:}

\begin{Shaded}
\begin{Highlighting}[]
\FunctionTok{sudo}\NormalTok{ apt install }\AttributeTok{{-}y}\NormalTok{ php php{-}cli php{-}mysql}
\ExtensionTok{php} \AttributeTok{{-}S}\NormalTok{ localhost:8000}
\end{Highlighting}
\end{Shaded}

\begin{center}\rule{0.5\linewidth}{0.5pt}\end{center}

\paragraph{7) IDE Setup}\label{ide-setup}

VS Code with Python/C++/ROS extensions for integrated linting,
debugging, and navigation.

\textbf{Install:} - VS Code: https://code.visualstudio.com/ - Browsers:
\texttt{sudo\ apt\ install\ -y\ firefox} - Arduino IDE:
\texttt{sudo\ apt\ install\ -y\ arduino}

\begin{center}\rule{0.5\linewidth}{0.5pt}\end{center}

\paragraph{8) Docker for Containerized
Workflows}\label{docker-for-containerized-workflows}

Use containers for reproducible environments in model serving and
simulation pipelines.

\begin{Shaded}
\begin{Highlighting}[]
\FunctionTok{sudo}\NormalTok{ apt install }\AttributeTok{{-}y}\NormalTok{ docker.io}
\FunctionTok{sudo}\NormalTok{ usermod }\AttributeTok{{-}aG}\NormalTok{ docker }\VariableTok{$USER}
\ExtensionTok{newgrp}\NormalTok{ docker}
\ExtensionTok{docker}\NormalTok{ run hello{-}world}
\ExtensionTok{docker}\NormalTok{ run }\AttributeTok{{-}it}\NormalTok{ ubuntu bash}
\end{Highlighting}
\end{Shaded}

\textbf{Distrobox (Run Any Distro Safely):}

\begin{Shaded}
\begin{Highlighting}[]
\FunctionTok{sudo}\NormalTok{ apt install }\AttributeTok{{-}y}\NormalTok{ podman}
\ExtensionTok{distrobox}\NormalTok{ create }\AttributeTok{{-}n}\NormalTok{ arch }\AttributeTok{{-}{-}image}\NormalTok{ archlinux:latest}
\ExtensionTok{distrobox}\NormalTok{ enter arch}
\end{Highlighting}
\end{Shaded}

Now you're in another distro. No VM. No pain.

\begin{center}\rule{0.5\linewidth}{0.5pt}\end{center}

\paragraph{9) Resource Optimization for Embedded
Targets}\label{resource-optimization-for-embedded-targets}

Account for CPU, RAM, thermal limits, and storage constraints before
deployment.

\begin{center}\rule{0.5\linewidth}{0.5pt}\end{center}

\paragraph{10) Update System \&
Repositories}\label{update-system-repositories}

\begin{quote}
→ \textbf{See Session 2 §4} for the full post-install checklist
including updates, codecs, Flatpak, Snap, and backups.
\end{quote}

\begin{Shaded}
\begin{Highlighting}[]
\FunctionTok{sudo}\NormalTok{ nano /etc/apt/sources.list    }\CommentTok{\# Edit sources}
\FunctionTok{sudo}\NormalTok{ apt update }\KeywordTok{\&\&} \FunctionTok{sudo}\NormalTok{ apt upgrade}
\end{Highlighting}
\end{Shaded}

\begin{quote}
\faExclamationCircle **Kali Repositories Warning:\textbf{ Don't mix Kali
repos with Ubuntu casually. You \emph{will} break things. Safe use case:
}specific tools only**. Use pinning if you \emph{must}.
\end{quote}

\begin{center}\rule{0.5\linewidth}{0.5pt}\end{center}

\subsubsection{Session 6: System Configuration \& Tools for
Robotics}\label{session-6-system-configuration-tools-for-robotics}

\paragraph{Learning outcomes}\label{learning-outcomes-5}

\begin{itemize}
\tightlist
\item
  Navigate Linux filesystem with robotics context
\item
  Manage permissions for hardware/device access
\item
  Configure shell environment and productivity shortcuts
\item
  Understand links and safe privilege usage
\end{itemize}

\begin{center}\rule{0.5\linewidth}{0.5pt}\end{center}

\paragraph{1) Filesystem Hierarchy in Robotics
Context}\label{filesystem-hierarchy-in-robotics-context}

\begin{quote}
→ \textbf{See Session 3} for hands-on navigation practice (\texttt{cd},
\texttt{ls}, \texttt{pwd}, paths). This section focuses on the
\emph{purpose} of each directory.
\end{quote}

Linux doesn't use drive letters like Windows
(\texttt{C:\textbackslash{}}, \texttt{D:\textbackslash{}}). Instead,
\textbf{everything starts at one place}: \texttt{/} (the root).

\textbf{\texttt{/} --- The Root of Everything}\\
\texttt{/} is the top of the file system tree. You \textbf{do not
casually mess around in \texttt{/}}.

\textbf{\texttt{/home} --- Where Humans Live}\\
This is where \textbf{users live}. Inside your home directory:
Documents, Downloads, Config files, SSH keys, Hidden dot files.
\textgreater{} • \textbf{Best Practice:} If you're learning Linux, live
in \texttt{/home}

\textbf{\texttt{/etc} --- Configuration Central}\\
Contains system configuration files, user settings, service configs,
network configs.

\begin{Shaded}
\begin{Highlighting}[]
\ExtensionTok{/etc/passwd}
\ExtensionTok{/etc/shadow}
\ExtensionTok{/etc/hosts}
\ExtensionTok{/etc/sudoers}
\end{Highlighting}
\end{Shaded}

\faExclamationCircle Editing files here can break login, networking, or
lock you out completely.\\
\textgreater{} • \textbf{Rule:} \texttt{/etc} = configuration, not
programs

\textbf{\texttt{/var} --- Variable Data (Stuff That Changes)}\\
Logs, caches, spools, databases.

\begin{Shaded}
\begin{Highlighting}[]
\ExtensionTok{/var/log}
\ExtensionTok{/var/log/syslog}
\ExtensionTok{/var/log/auth.log}
\end{Highlighting}
\end{Shaded}

If a system is acting weird → \textbf{Check \texttt{/var/log} or related
logs}

\textbf{\texttt{/usr} --- User Programs (Not Users)}\\
Confusing name---this is \textbf{not user home directories}. Contains
installed programs, libraries, shared resources.

\begin{Shaded}
\begin{Highlighting}[]
\ExtensionTok{/usr/bin}
\ExtensionTok{/usr/lib}
\ExtensionTok{/usr/share}
\end{Highlighting}
\end{Shaded}

\textbf{\texttt{/dev/} --- Hardware/Device Nodes} (serial, camera,
input)\\
\textbf{\texttt{/sys/class/gpio/} --- GPIO representation} (conceptual
Pi mapping)\\
\textbf{\texttt{/opt/ros/} --- ROS installation location}

\begin{center}\rule{0.5\linewidth}{0.5pt}\end{center}

\paragraph{2) ``Everything Is a File'' (Linux
Philosophy)}\label{everything-is-a-file-linux-philosophy}

\begin{quote}
→ \textbf{See Session 3 §10} for the full explanation with
\texttt{/dev/}, \texttt{/proc/}, and \texttt{/sys/} examples.
\end{quote}

This is where Linux becomes \emph{cool}. In Linux, files, directories,
devices, hardware, and network interfaces are all represented as files.

\begin{Shaded}
\begin{Highlighting}[]
\ExtensionTok{/dev/sda}          \CommentTok{\# Hard drive}
\ExtensionTok{/dev/null}         \CommentTok{\# Black hole}
\ExtensionTok{/proc/cpuinfo}     \CommentTok{\# CPU info}
\ExtensionTok{/dev/tty}          \CommentTok{\# Current terminal}
\ExtensionTok{/dev/random}       \CommentTok{\# Random number generator}
\ExtensionTok{/dev/urandom}      \CommentTok{\# Faster random numbers}
\ExtensionTok{/proc/uptime}      \CommentTok{\# System uptime}
\ExtensionTok{/proc/net/tcp}     \CommentTok{\# TCP connections}
\ExtensionTok{/sys/class/net/eth0/address}         \CommentTok{\# Network MAC}
\ExtensionTok{/sys/class/backlight/brightness}     \CommentTok{\# Screen brightness}
\ExtensionTok{/sys/class/power\_supply/BAT0/capacity} \CommentTok{\# Battery level}
\ExtensionTok{/dev/loop0}        \CommentTok{\# Loopback device (mount files as disks)}
\end{Highlighting}
\end{Shaded}

This design is WHY Linux is scriptable, automatable, and runs the
internet.

\begin{center}\rule{0.5\linewidth}{0.5pt}\end{center}

\paragraph{3) Dotfiles (Hidden Power)}\label{dotfiles-hidden-power}

\begin{quote}
→ \textbf{See Session 3 §12} where \texttt{ls\ -la} first introduces
hidden files. This section goes deeper.
\end{quote}

Files that start with a dot \texttt{.} are \textbf{hidden}.

\begin{Shaded}
\begin{Highlighting}[]
\ExtensionTok{.bashrc}
\ExtensionTok{.profile}
\ExtensionTok{.gitconfig}
\ExtensionTok{.ssh/}
\end{Highlighting}
\end{Shaded}

They usually store shell configs, app preferences, and user settings.
List them: \texttt{ls\ -a}

Why hide them? Prevent clutter and accidental deletion.\\
\textgreater{} • \textbf{Important:} Deleting dot files can reset your
environment or break tools.

\textbf{Dotfiles (\texttt{.bashrc}, \texttt{.profile}):} Used to
customize shell behavior, aliases, environment variables, and startup
logic.

\begin{center}\rule{0.5\linewidth}{0.5pt}\end{center}

\paragraph{4) Permissions and Groups for Hardware
Access}\label{permissions-and-groups-for-hardware-access}

\textbf{Why This Matters (Seriously):}\\
Linux doesn't break because it's fragile. Linux breaks because
\textbf{permissions exist} and \textbf{you don't understand them yet}.

Permissions protect your system, prevent malware from wrecking
everything, and make Linux powerful for servers, DevOps, security, and
cloud. Once this clicks: errors make sense, ``Permission denied'' stops
ruining your day, and you stop nuking your system with
\texttt{sudo\ rm\ -rf}.

Every file has: - \textbf{Owner} - \textbf{Group} - \textbf{Permissions}

Check with:

\begin{Shaded}
\begin{Highlighting}[]
\FunctionTok{ls} \AttributeTok{{-}l}
\end{Highlighting}
\end{Shaded}

Example:

\begin{Shaded}
\begin{Highlighting}[]
\ExtensionTok{{-}rwxr{-}xr{-}{-}}\NormalTok{ 1 ahmad staff script.sh}
\end{Highlighting}
\end{Shaded}

Breakdown: - Owner: ahmad (rwx) - Group: staff (r-x) - Others: (r--)

\textbf{Three permission types:} - \textbf{r} = read (4) - \textbf{w} =
write (2) - \textbf{x} = execute (1)

\textbf{Permission Numbers (Octal Mode):}

\begin{verbatim}
7 = rwx
6 = rw-
5 = r-x
4 = r--
\end{verbatim}

\begin{Shaded}
\begin{Highlighting}[]
\FunctionTok{chmod}\NormalTok{ 755 script.sh    }\CommentTok{\# Owner: rwx, Group: r{-}x, Others: r{-}x}
\FunctionTok{chmod}\NormalTok{ +x script.sh     }\CommentTok{\# Make executable}
\end{Highlighting}
\end{Shaded}

\begin{quote}
• You will see \texttt{755} and \texttt{644} everywhere.
\end{quote}

\textbf{Changing Ownership:}

\begin{Shaded}
\begin{Highlighting}[]
\FunctionTok{chown}\NormalTok{ user file}
\FunctionTok{chown}\NormalTok{ user:group file}
\FunctionTok{sudo}\NormalTok{ chown root:root config.conf}
\end{Highlighting}
\end{Shaded}

\faExclamationCircle Wrong ownership = broken apps

\textbf{Typical groups for hardware access:} \texttt{dialout},
\texttt{gpio}, \texttt{video}. Correct group membership often resolves
``permission denied'' with serial/camera devices.

\begin{center}\rule{0.5\linewidth}{0.5pt}\end{center}

\paragraph{\texorpdfstring{5) Root Permissions (\texttt{sudo},
\texttt{su},
\texttt{/etc/sudoers})}{5) Root Permissions (sudo, su, /etc/sudoers)}}\label{root-permissions-sudo-su-etcsudoers}

\textbf{Root} is the system administrator---the power mode user. Root
can read any file, delete any file, and change any permission.

\textbf{\texttt{sudo} --- Temporary Root:}

\begin{Shaded}
\begin{Highlighting}[]
\FunctionTok{sudo}\NormalTok{ command}
\end{Highlighting}
\end{Shaded}

Executes command as root, logs actions, requires password.\\
\textgreater{} • \textbf{Best practice:} Use \texttt{sudo}, not
permanent root login

\textbf{\texttt{su} --- Switch User:}

\begin{Shaded}
\begin{Highlighting}[]
\FunctionTok{su}
\FunctionTok{su}\NormalTok{ username}
\end{Highlighting}
\end{Shaded}

Switches shell to another user. \faExclamationCircle Dangerous if
misused.

\textbf{\texttt{/etc/sudoers} --- Who Can Use sudo:}\\
DO NOT edit directly. Always use:

\begin{Shaded}
\begin{Highlighting}[]
\FunctionTok{sudo}\NormalTok{ visudo}
\end{Highlighting}
\end{Shaded}

Wrong syntax here → no sudo → bad day.

\textbf{Common Beginner Mistakes:} - \faTimesCircle Running everything
with sudo (bad for security) - \faTimesCircle \texttt{chmod\ 777}
everything (everyone can do everything---bad for security \& interviews
:-)) - \faTimesCircle Editing system files without backup
(\texttt{sudo\ cp\ config.conf\ config.conf.bak}---always!) -
\faTimesCircle Deleting random stuff in \texttt{/etc} or \texttt{/usr} -
\faTimesCircle Not reading error messages (Linux tells you what failed
and why)

\begin{quote}
• Use sudo \textbf{only when required}. Permissions are not
obstacles---they are \textbf{guardrails}.
\end{quote}

\begin{center}\rule{0.5\linewidth}{0.5pt}\end{center}

\paragraph{6) Hard Links and Soft
Links}\label{hard-links-and-soft-links}

Links are \textbf{references} to files.

\textbf{Hard Links:} - Point to the same inode - File exists in multiple
locations - Deleting one doesn't remove data

\begin{Shaded}
\begin{Highlighting}[]
\FunctionTok{ln}\NormalTok{ file1 file2}
\end{Highlighting}
\end{Shaded}

Limitations: Same filesystem only, can't link directories.

\textbf{Soft Links (Symlinks):} - Pointer to file path - Like Windows
shortcuts

\begin{Shaded}
\begin{Highlighting}[]
\FunctionTok{ln} \AttributeTok{{-}s}\NormalTok{ target linkname}
\FunctionTok{ln} \AttributeTok{{-}s}\NormalTok{ /var/log/syslog syslog\_link}
\end{Highlighting}
\end{Shaded}

If target is deleted → \faTimesCircle Link breaks\\
\textgreater{} • Most commonly used type.

\begin{center}\rule{0.5\linewidth}{0.5pt}\end{center}

\paragraph{7) Shell Aliases}\label{shell-aliases}

\begin{quote}
→ \textbf{See Session 1 §11} for the mindset behind customization and
efficiency.
\end{quote}

Aliases save time---like macros for your shell. Speed up common
workflows with carefully chosen aliases.

\begin{Shaded}
\begin{Highlighting}[]
\BuiltInTok{alias}\NormalTok{ ll=}\StringTok{\textquotesingle{}ls {-}la\textquotesingle{}}
\BuiltInTok{alias}\NormalTok{ gs=}\StringTok{\textquotesingle{}git status\textquotesingle{}}
\BuiltInTok{echo} \StringTok{"alias ll=\textquotesingle{}ls {-}lah\textquotesingle{}"} \OperatorTok{\textgreater{}\textgreater{}}\NormalTok{ \textasciitilde{}/.bashrc}
\end{Highlighting}
\end{Shaded}

Add to \texttt{\textasciitilde{}/.bashrc} or
\texttt{\textasciitilde{}/.zshrc} and
\texttt{source\ \textasciitilde{}/.bashrc} to reload. You'll thank
yourself later when typing gets lazy :)

\begin{center}\rule{0.5\linewidth}{0.5pt}\end{center}

\paragraph{8) Font/Theme Customization}\label{fonttheme-customization}

Long robotics sessions benefit from readable fonts and reduced visual
fatigue.

\begin{itemize}
\tightlist
\item
  Install custom fonts: \texttt{sudo\ apt\ install\ fonts-roboto}
\item
  GNOME Tweaks for themes
\item
  Flatpak themes from Flathub
\end{itemize}

Fonts + themes = aesthetic \textasciitilde100\% more hipster points B-)

\begin{center}\rule{0.5\linewidth}{0.5pt}\end{center}

\paragraph{9) Virtual Device
Exploration}\label{virtual-device-exploration}

Understand how simulated hardware appears in \texttt{/dev} and how
tooling discovers it.

\begin{center}\rule{0.5\linewidth}{0.5pt}\end{center}

\paragraph{10) Install Common Software}\label{install-common-software}

\begin{quote}
→ \textbf{See Session 3 §7} for detailed package manager coverage
(\texttt{apt}, \texttt{dnf}, \texttt{yum}, \texttt{pip},
\texttt{snap}).\\
→ \textbf{See Session 2 §4} for Flatpak and Snap setup details.
\end{quote}

\textbf{Quick starter pack for Ubuntu:}

\begin{Shaded}
\begin{Highlighting}[]
\FunctionTok{sudo}\NormalTok{ apt update}
\FunctionTok{sudo}\NormalTok{ apt install }\DataTypeTok{\textbackslash{}}
\NormalTok{    git htop tmux fastfetch}
\end{Highlighting}
\end{Shaded}

\begin{center}\rule{0.5\linewidth}{0.5pt}\end{center}

\paragraph{Useful Links --- Session 6}\label{useful-links-session-6}

\textbf{Package Management:}

%
  \par\vspace{0.5ex}%
  \href{https://wiki.debian.org/Apt}{\textcolor{blue}{\underline{\faLink\quad APT docs (Debian Wiki)}}}%
  \par\vspace{0.5ex}%

%
  \par\vspace{0.5ex}%
  \href{https://help.ubuntu.com/community/AptGet/Howto}{\textcolor{blue}{\underline{\faLink\quad APT Howto (Ubuntu Community)}}}%
  \par\vspace{0.5ex}%

%
  \par\vspace{0.5ex}%
  \href{https://wiki.debian.org/SourcesList}{\textcolor{blue}{\underline{\faLink\quad Debian SourcesList}}}%
  \par\vspace{0.5ex}%

%
  \par\vspace{0.5ex}%
  \href{https://help.ubuntu.com/community/Repositories}{\textcolor{blue}{\underline{\faLink\quad Ubuntu Repositories}}}%
  \par\vspace{0.5ex}%

%
  \par\vspace{0.5ex}%
  \href{https://www.kali.org/docs/general-use/kali-linux-sources-list-repositories/}{\textcolor{blue}{\underline{\faLink\quad Kali Repositories docs}}}%
  \par\vspace{0.5ex}%

%
  \par\vspace{0.5ex}%
  \href{https://youtu.be/YI1Q3R0TEYs?si=TO3KfOz4nq0sSIVw}{\textcolor{blue}{\underline{\faLink\quad Adding Kali repos to Debian (YouTube)}}}%
  \par\vspace{0.5ex}%

\textbf{Build from Source:}

%
  \par\vspace{0.5ex}%
  \href{https://suckless.org/st/}{\textcolor{blue}{\underline{\faLink\quad suckless st}}}%
  \par\vspace{0.5ex}%

%
  \par\vspace{0.5ex}%
  \href{https://suckless.org/dmenu/}{\textcolor{blue}{\underline{\faLink\quad suckless dmenu}}}%
  \par\vspace{0.5ex}%

%
  \par\vspace{0.5ex}%
  \href{https://gcc.gnu.org/onlinedocs/}{\textcolor{blue}{\underline{\faLink\quad GCC docs}}}%
  \par\vspace{0.5ex}%

%
  \par\vspace{0.5ex}%
  \href{https://www.gnu.org/software/gdb/documentation/}{\textcolor{blue}{\underline{\faLink\quad GDB docs}}}%
  \par\vspace{0.5ex}%

%
  \par\vspace{0.5ex}%
  \href{https://www.gnu.org/software/make/manual/}{\textcolor{blue}{\underline{\faLink\quad GNU Make manual}}}%
  \par\vspace{0.5ex}%

%
  \par\vspace{0.5ex}%
  \href{https://www.nongnu.org/avr-libc/}{\textcolor{blue}{\underline{\faLink\quad avr-libc}}}%
  \par\vspace{0.5ex}%

\textbf{Web Development:}

%
  \par\vspace{0.5ex}%
  \href{https://docs.python.org/3/library/http.server.html}{\textcolor{blue}{\underline{\faLink\quad Python http.server docs}}}%
  \par\vspace{0.5ex}%

%
  \par\vspace{0.5ex}%
  \href{https://www.npmjs.com/package/serve}{\textcolor{blue}{\underline{\faLink\quad npm serve}}}%
  \par\vspace{0.5ex}%

%
  \par\vspace{0.5ex}%
  \href{https://www.php.net/manual/en/}{\textcolor{blue}{\underline{\faLink\quad PHP Manual}}}%
  \par\vspace{0.5ex}%

%
  \par\vspace{0.5ex}%
  \href{https://www.apachefriends.org/index.html}{\textcolor{blue}{\underline{\faLink\quad XAMPP download}}}%
  \par\vspace{0.5ex}%

%
  \par\vspace{0.5ex}%
  \href{https://www.apachefriends.org/faq\_linux.html}{\textcolor{blue}{\underline{\faLink\quad XAMPP FAQ (Linux)}}}%
  \par\vspace{0.5ex}%

\textbf{Node.js \& React:}

%
  \par\vspace{0.5ex}%
  \href{https://nodejs.org/en/docs}{\textcolor{blue}{\underline{\faLink\quad Node.js docs}}}%
  \par\vspace{0.5ex}%

%
  \par\vspace{0.5ex}%
  \href{https://react.dev/learn}{\textcolor{blue}{\underline{\faLink\quad React docs}}}%
  \par\vspace{0.5ex}%

\textbf{Java:}

%
  \par\vspace{0.5ex}%
  \href{https://openjdk.org/}{\textcolor{blue}{\underline{\faLink\quad OpenJDK}}}%
  \par\vspace{0.5ex}%

%
  \par\vspace{0.5ex}%
  \href{https://docs.oracle.com/javase/tutorial/}{\textcolor{blue}{\underline{\faLink\quad Oracle Java Tutorial}}}%
  \par\vspace{0.5ex}%

\textbf{Python:}

%
  \par\vspace{0.5ex}%
  \href{https://docs.python.org/3/tutorial/venv.html}{\textcolor{blue}{\underline{\faLink\quad Python venv tutorial}}}%
  \par\vspace{0.5ex}%

%
  \par\vspace{0.5ex}%
  \href{https://pip.pypa.io/en/stable/}{\textcolor{blue}{\underline{\faLink\quad pip docs}}}%
  \par\vspace{0.5ex}%

\textbf{Apps \& IDEs:}

%
  \par\vspace{0.5ex}%
  \href{https://www.google.com/chrome/}{\textcolor{blue}{\underline{\faLink\quad Google Chrome}}}%
  \par\vspace{0.5ex}%

%
  \par\vspace{0.5ex}%
  \href{https://code.visualstudio.com/}{\textcolor{blue}{\underline{\faLink\quad VS Code}}}%
  \par\vspace{0.5ex}%

%
  \par\vspace{0.5ex}%
  \href{https://code.visualstudio.com/docs/setup/linux}{\textcolor{blue}{\underline{\faLink\quad VS Code Linux Setup}}}%
  \par\vspace{0.5ex}%

%
  \par\vspace{0.5ex}%
  \href{https://docs.arduino.cc/}{\textcolor{blue}{\underline{\faLink\quad Arduino}}}%
  \par\vspace{0.5ex}%

\textbf{Containers:}

%
  \par\vspace{0.5ex}%
  \href{https://docs.docker.com/engine/install/ubuntu/}{\textcolor{blue}{\underline{\faLink\quad Docker install docs}}}%
  \par\vspace{0.5ex}%

%
  \par\vspace{0.5ex}%
  \href{https://mmbesar.github.io/Tutorials/HS-Docker/}{\textcolor{blue}{\underline{\faLink\quad Docker tutorial (mmbesar)}}}%
  \par\vspace{0.5ex}%

%
  \par\vspace{0.5ex}%
  \href{https://distrobox.it/}{\textcolor{blue}{\underline{\faLink\quad Distrobox docs}}}%
  \par\vspace{0.5ex}%

\textbf{AppImage:}

%
  \par\vspace{0.5ex}%
  \href{https://docs.appimage.org/}{\textcolor{blue}{\underline{\faLink\quad AppImage docs}}}%
  \par\vspace{0.5ex}%

%
  \par\vspace{0.5ex}%
  \href{https://mmbesar.github.io/Tutorials/appimages/}{\textcolor{blue}{\underline{\faLink\quad AppImage tutorial (mmbesar)}}}%
  \par\vspace{0.5ex}%

\begin{center}\rule{0.5\linewidth}{0.5pt}\end{center}

\subsubsection{Session 7: Process Monitoring \& Robotics System
Management}\label{session-7-process-monitoring-robotics-system-management}

\begin{quote}
\textbf{Verify \texttt{htop} is installed.}
\end{quote}

\paragraph{Learning outcomes}\label{learning-outcomes-6}

\begin{itemize}
\tightlist
\item
  Monitor process/resource usage under robotics workloads
\item
  Control or terminate stalled tasks safely
\item
  Use background execution and command chaining
\item
  Reason about embedded constraints (thermal/memory throttling)
\end{itemize}

\begin{center}\rule{0.5\linewidth}{0.5pt}\end{center}

\paragraph{1) What Is a Process?}\label{what-is-a-process}

Every program you run becomes a \emph{process}---an instance of a
running program with a unique ID (PID). Linux lets you \textbf{list},
\textbf{monitor}, and \textbf{control} these processes using simple
tools.

\begin{center}\rule{0.5\linewidth}{0.5pt}\end{center}

\paragraph{\texorpdfstring{2) \texttt{htop} --- Interactive Process
Viewer}{2) htop --- Interactive Process Viewer}}\label{htop-interactive-process-viewer}

\faLightbulb \texttt{htop} is like \texttt{top}, but \emph{way more
human-friendly}: color, scrolling, arrow navigation, signal menus, and
killer shortcuts.

\textbf{Why \texttt{htop}?} - Shows CPU, RAM, threads in real-time -
Scroll through processes - Select and kill without typing PIDs

\textbf{Install \& Run:}

\begin{Shaded}
\begin{Highlighting}[]
\FunctionTok{sudo}\NormalTok{ apt update}
\FunctionTok{sudo}\NormalTok{ apt install htop}
\ExtensionTok{htop}
\end{Highlighting}
\end{Shaded}

\textbf{Inside \texttt{htop}:} - ↑/↓ --- move - F9 --- kill selected
process - F3 --- search/filter - F5 --- tree view (parent/child
structure)

\begin{quote}
Always use \texttt{htop} to \textbf{see what's happening right
now}---you'll \emph{instantly} understand resource bottlenecks.
\end{quote}

\begin{center}\rule{0.5\linewidth}{0.5pt}\end{center}

\paragraph{\texorpdfstring{3) \texttt{kill} and \texttt{pkill} --- End
Processes}{3) kill and pkill --- End Processes}}\label{kill-and-pkill-end-processes}

When a process misbehaves (freezes, uses all your CPU), you kill it.

\textbf{kill} (targets by PID):

\begin{Shaded}
\begin{Highlighting}[]
\BuiltInTok{kill}\NormalTok{ 1234               }\CommentTok{\# SIGTERM {-} politely asks to quit}
\BuiltInTok{kill} \AttributeTok{{-}9}\NormalTok{ 1234            }\CommentTok{\# SIGKILL {-} stops immediately}
\end{Highlighting}
\end{Shaded}

\textbf{pkill} (targets by name .. no need to look up PID):

\begin{Shaded}
\begin{Highlighting}[]
\ExtensionTok{pkill}\NormalTok{ firefox}
\ExtensionTok{pkill} \AttributeTok{{-}u}\NormalTok{ someuser}
\end{Highlighting}
\end{Shaded}

{[}NOTE{]} Tip: \texttt{kill} works on PIDs, \texttt{pkill} works on
\emph{names}.\\
\textgreater{} \faExclamationCircle Use \texttt{kill}
sparingly---force-killing can corrupt data if the program is mid-write.

\begin{center}\rule{0.5\linewidth}{0.5pt}\end{center}

\paragraph{\texorpdfstring{4) \texttt{ps} + \texttt{grep} --- Find
Processes via
CLI}{4) ps + grep --- Find Processes via CLI}}\label{ps-grep-find-processes-via-cli}

\texttt{ps} gives you a \emph{snapshot} of currently running
processes---not real-time like \texttt{htop}, but extremely useful.

\begin{Shaded}
\begin{Highlighting}[]
\FunctionTok{ps}\NormalTok{ aux                    }\CommentTok{\# List ALL processes}
\FunctionTok{ps}\NormalTok{ aux }\KeywordTok{|} \FunctionTok{grep} \AttributeTok{{-}i}\NormalTok{ ros      }\CommentTok{\# Find ROS processes}
\FunctionTok{ps}\NormalTok{ aux }\KeywordTok{|} \FunctionTok{grep}\NormalTok{ firefox     }\CommentTok{\# Find Firefox}
\FunctionTok{ps} \AttributeTok{{-}ef} \KeywordTok{|} \FunctionTok{grep}\NormalTok{ python      }\CommentTok{\# Find Python instances}
\end{Highlighting}
\end{Shaded}

\faExclamationCircle Pro tip: wrap your grep to \emph{not} match itself
(e.g., \texttt{grep\ "{[}f{]}irefox"}), because normal grep often shows
its \emph{own} process.

\begin{center}\rule{0.5\linewidth}{0.5pt}\end{center}

\paragraph{\texorpdfstring{5) Background Execution
(\texttt{\&})}{5) Background Execution (\&)}}\label{background-execution}

Sometimes you want to start a long-running task \emph{without blocking
your terminal}.

\begin{Shaded}
\begin{Highlighting}[]
\FunctionTok{sleep}\NormalTok{ 60 }\KeywordTok{\&}                \CommentTok{\# Run in background}
\end{Highlighting}
\end{Shaded}

\begin{itemize}
\tightlist
\item
  \texttt{\&} → run command in the background
\item
  You get your prompt back immediately
\item
  Shell assigns a job number and PID
\end{itemize}

\textbf{Job Management:}

\begin{Shaded}
\begin{Highlighting}[]
\BuiltInTok{jobs}                      \CommentTok{\# List jobs}
\BuiltInTok{fg}\NormalTok{ \%1                     }\CommentTok{\# Bring job 1 to foreground}
\BuiltInTok{bg}\NormalTok{ \%1                     }\CommentTok{\# Resume job 1 in background}
\end{Highlighting}
\end{Shaded}

\texttt{\&}, \texttt{screen}, \texttt{tmux} keep long tasks active while
terminal remains usable.

\begin{center}\rule{0.5\linewidth}{0.5pt}\end{center}

\paragraph{\texorpdfstring{6) Command Chaining --- \texttt{\&\&},
\texttt{\textbar{}\textbar{}},
\texttt{;}}{6) Command Chaining --- \&\&, \textbar\textbar, ;}}\label{command-chaining}

Command chaining lets you fire off multiple commands in one
line---\emph{with logic about success or failure}.

{\def\LTcaptype{none} % do not increment counter
\begin{longtable}[]{@{}ll@{}}
\toprule\noalign{}
Operator & What it Means \\
\midrule\noalign{}
\endhead
\bottomrule\noalign{}
\endlastfoot
\texttt{cmd1\ ;\ cmd2} & Run cmd1, then cmd2 no matter what \\
\texttt{cmd1\ \&\&\ cmd2} & Run cmd2 ONLY if cmd1 succeeds (exit status
0) \\
\texttt{cmd1\ \textbackslash{}\textbar{}\textbackslash{}\textbar{}\ cmd2}
& Run cmd2 ONLY if cmd1 fails \\
\end{longtable}
}

\begin{Shaded}
\begin{Highlighting}[]
\FunctionTok{sudo}\NormalTok{ apt update }\KeywordTok{\&\&} \FunctionTok{sudo}\NormalTok{ apt upgrade }\AttributeTok{{-}y}
\FunctionTok{rm}\NormalTok{ temp.log }\KeywordTok{||} \BuiltInTok{echo} \StringTok{"Could not delete temp.log"}
\ExtensionTok{cmd1} \KeywordTok{;} \ExtensionTok{cmd2} \KeywordTok{;} \ExtensionTok{cmd3}
\end{Highlighting}
\end{Shaded}

• Perfect for chaining updates, scripts, and error handling. Practice
chaining---it's a \emph{game changer} for automation and scripting.

\begin{center}\rule{0.5\linewidth}{0.5pt}\end{center}

\paragraph{7) Embedded Constraints}\label{embedded-constraints}

Understand temperature, throttling, and memory pressure behavior on
small devices.

\begin{center}\rule{0.5\linewidth}{0.5pt}\end{center}

\paragraph{8) Real-Time Resource
Management}\label{real-time-resource-management}

Prioritize critical processes and avoid overload from concurrent jobs.

\begin{center}\rule{0.5\linewidth}{0.5pt}\end{center}

\paragraph{9) Simulation Monitoring}\label{simulation-monitoring}

Track simulator CPU/RAM impact and tune settings to maintain
responsiveness.

\begin{center}\rule{0.5\linewidth}{0.5pt}\end{center}

\paragraph{Useful Links --- Session 7}\label{useful-links-session-7}

\textbf{Process Management:}

%
  \par\vspace{0.5ex}%
  \href{https://www.geeksforgeeks.org/linux-unix/using-htop-to-monitor-system-processes-on-linux/}{\textcolor{blue}{\underline{\faLink\quad htop tutorial (GeeksforGeeks)}}}%
  \par\vspace{0.5ex}%

%
  \par\vspace{0.5ex}%
  \href{https://www.howtogeek.com/how-to-manage-linux-processes-using-ps-kill-and-pkill/}{\textcolor{blue}{\underline{\faLink\quad Manage processes with ps, kill, pkill (How-To Geek)}}}%
  \par\vspace{0.5ex}%

%
  \par\vspace{0.5ex}%
  \href{https://discourse.ubuntu.com/t/viewing-and-monitoring-processes-in-linux/26024}{\textcolor{blue}{\underline{\faLink\quad Viewing and Monitoring Processes (Ubuntu Community)}}}%
  \par\vspace{0.5ex}%

\textbf{Job Control \& Commands:}

%
  \par\vspace{0.5ex}%
  \href{https://www.digitalocean.com/community/tutorials/how-to-use-bash-s-job-control-to-manage-foreground-and-background-processes}{\textcolor{blue}{\underline{\faLink\quad Bash job control (DigitalOcean)}}}%
  \par\vspace{0.5ex}%

%
  \par\vspace{0.5ex}%
  \href{https://www.geeksforgeeks.org/linux-unix/chaining-commands-in-linux/}{\textcolor{blue}{\underline{\faLink\quad Chaining Commands (GeeksforGeeks)}}}%
  \par\vspace{0.5ex}%

\textbf{General References:}

%
  \par\vspace{0.5ex}%
  \href{https://ubuntu.com/tutorials/command-line-for-beginners}{\textcolor{blue}{\underline{\faLink\quad Ubuntu CLI for Beginners}}}%
  \par\vspace{0.5ex}%

%
  \par\vspace{0.5ex}%
  \href{https://www.gnu.org/software/bash/manual/}{\textcolor{blue}{\underline{\faLink\quad Bash Reference Manual}}}%
  \par\vspace{0.5ex}%

\begin{center}\rule{0.5\linewidth}{0.5pt}\end{center}

\subsection{Chapter 3: System Management, Scripting \&
Networking}\label{chapter-3-system-management-scripting-networking}

\subsubsection{Session 8: Automation with systemd \&
Cron}\label{session-8-automation-with-systemd-cron}

\begin{quote}
Alright .. this is the session where you level up from ``I run Linux''
to: ``I control what runs on Linux\ldots{} and when it runs.''
\end{quote}

\paragraph{Learning outcomes}\label{learning-outcomes-7}

\begin{itemize}
\tightlist
\item
  Create/manage systemd services for robotics startup
\item
  Schedule recurring tasks using timers and cron
\item
  Understand unit dependencies and debugging basics
\end{itemize}

\begin{center}\rule{0.5\linewidth}{0.5pt}\end{center}

\paragraph{1) What is systemd? (And Why You Should
Care)}\label{what-is-systemd-and-why-you-should-care}

On modern Debian-based systems (Debian, Ubuntu, Linux Mint), the init
system is \textbf{systemd}.

An \emph{init system} is the first userspace process started by the
kernel. Check it:

\begin{Shaded}
\begin{Highlighting}[]
\FunctionTok{ps} \AttributeTok{{-}p}\NormalTok{ 1 }\AttributeTok{{-}o}\NormalTok{ comm=    }\CommentTok{\# Should show "systemd"}
\end{Highlighting}
\end{Shaded}

\textbf{What Does systemd Do?} - Starts services during boot - Manages
background processes (daemons) - Handles logging (via \texttt{journald})
- Mounts filesystems - Manages network services - Handles timers (our
main topic today)

\begin{quote}
If Linux is a city □, \texttt{systemd} is the mayor.
\end{quote}

\begin{center}\rule{0.5\linewidth}{0.5pt}\end{center}

\paragraph{2) Understanding Units in
systemd}\label{understanding-units-in-systemd}

Everything in systemd is a \textbf{Unit}---a configuration file that
tells systemd what to manage.

{\def\LTcaptype{none} % do not increment counter
\begin{longtable}[]{@{}lll@{}}
\toprule\noalign{}
Type & Extension & Purpose \\
\midrule\noalign{}
\endhead
\bottomrule\noalign{}
\endlastfoot
Service & \texttt{.service} & Background services \\
Timer & \texttt{.timer} & Scheduling \\
Mount & \texttt{.mount} & Mount points \\
Target & \texttt{.target} & Group of units \\
Socket & \texttt{.socket} & Socket activation \\
\end{longtable}
}

\begin{Shaded}
\begin{Highlighting}[]
\ExtensionTok{systemctl}\NormalTok{ list{-}units           }\CommentTok{\# List all active units}
\ExtensionTok{systemctl}\NormalTok{ list{-}unit{-}files      }\CommentTok{\# List all installed units}
\end{Highlighting}
\end{Shaded}

\begin{center}\rule{0.5\linewidth}{0.5pt}\end{center}

\paragraph{3) Managing Services with
systemctl}\label{managing-services-with-systemctl}

This is daily Linux admin stuff.

\begin{Shaded}
\begin{Highlighting}[]
\FunctionTok{sudo}\NormalTok{ systemctl start apache2       }\CommentTok{\# Start}
\FunctionTok{sudo}\NormalTok{ systemctl stop apache2        }\CommentTok{\# Stop}
\FunctionTok{sudo}\NormalTok{ systemctl restart apache2     }\CommentTok{\# Restart}
\FunctionTok{sudo}\NormalTok{ systemctl reload apache2      }\CommentTok{\# Reload config (keeps connections alive)}
\FunctionTok{sudo}\NormalTok{ systemctl enable apache2      }\CommentTok{\# Enable at boot}
\FunctionTok{sudo}\NormalTok{ systemctl disable apache2     }\CommentTok{\# Disable at boot}
\FunctionTok{sudo}\NormalTok{ systemctl mask apache2        }\CommentTok{\# Completely block from starting}
\FunctionTok{sudo}\NormalTok{ systemctl unmask apache2      }\CommentTok{\# Unmask}
\ExtensionTok{systemctl}\NormalTok{ status apache2           }\CommentTok{\# Check status}
\end{Highlighting}
\end{Shaded}

\textbf{Reload vs Restart (Important!):} - \texttt{reload}: Reloads
configuration, does NOT stop process, keeps connections alive -
\texttt{restart}: Stops service, starts again, drops active connections

\begin{center}\rule{0.5\linewidth}{0.5pt}\end{center}

\paragraph{4) Where Are Unit Files
Stored?}\label{where-are-unit-files-stored}

{\def\LTcaptype{none} % do not increment counter
\begin{longtable}[]{@{}ll@{}}
\toprule\noalign{}
Directory & Purpose \\
\midrule\noalign{}
\endhead
\bottomrule\noalign{}
\endlastfoot
\texttt{/usr/lib/systemd/system/} & Default package units \\
\texttt{/lib/systemd/system/} & (Debian equivalent) \\
\texttt{/etc/systemd/system/} & Custom \& overrides \\
\end{longtable}
}

\textbf{Priority Order (highest to lowest):}
\texttt{/etc/systemd/system/} overrides defaults.

\begin{center}\rule{0.5\linewidth}{0.5pt}\end{center}

\paragraph{5) Anatomy of a Service
File}\label{anatomy-of-a-service-file}

Inspect a service file structure:

\begin{Shaded}
\begin{Highlighting}[]
\ExtensionTok{systemctl}\NormalTok{ cat apache2.service}
\end{Highlighting}
\end{Shaded}

Typical structure:

\begin{Shaded}
\begin{Highlighting}[]
\KeywordTok{[Unit]}
\DataTypeTok{Description}\OtherTok{=}\StringTok{...}
\DataTypeTok{After}\OtherTok{=}\StringTok{network.target}

\KeywordTok{[Service]}
\DataTypeTok{Type}\OtherTok{=}\StringTok{forking}
\DataTypeTok{ExecStart}\OtherTok{=}\StringTok{/usr/sbin/apachectl start}
\DataTypeTok{ExecReload}\OtherTok{=}\StringTok{/usr/sbin/apachectl graceful}
\DataTypeTok{ExecStop}\OtherTok{=}\StringTok{/usr/sbin/apachectl stop}

\KeywordTok{[Install]}
\DataTypeTok{WantedBy}\OtherTok{=}\StringTok{multi{-}user.target}
\end{Highlighting}
\end{Shaded}

\textbf{{[}Unit{]} Section:} Metadata and dependencies (Description,
After=, Requires=)\\
\textbf{{[}Service{]} Section:} How the service runs (ExecStart,
ExecStop, Restart, User, Type)\\
\textbf{{[}Install{]} Section:} Boot integration
(\texttt{WantedBy=multi-user.target} means start in normal multi-user
mode)

\begin{center}\rule{0.5\linewidth}{0.5pt}\end{center}

\paragraph{6) Creating Your Own
Service}\label{creating-your-own-service}

\textbf{Goal: Run a simple script at boot}

Create script:

\begin{Shaded}
\begin{Highlighting}[]
\FunctionTok{sudo}\NormalTok{ nano /usr/local/bin/hello.sh}
\end{Highlighting}
\end{Shaded}

\begin{Shaded}
\begin{Highlighting}[]
\CommentTok{\#!/bin/bash}
\BuiltInTok{echo} \StringTok{"Hello from systemd!"} \OperatorTok{\textgreater{}\textgreater{}}\NormalTok{ /tmp/hello.log}
\end{Highlighting}
\end{Shaded}

Make executable:

\begin{Shaded}
\begin{Highlighting}[]
\FunctionTok{sudo}\NormalTok{ chmod +x /usr/local/bin/hello.sh}
\end{Highlighting}
\end{Shaded}

Create service file:

\begin{Shaded}
\begin{Highlighting}[]
\FunctionTok{sudo}\NormalTok{ nano /etc/systemd/system/hello.service}
\end{Highlighting}
\end{Shaded}

\begin{Shaded}
\begin{Highlighting}[]
\KeywordTok{[Unit]}
\DataTypeTok{Description}\OtherTok{=}\StringTok{My Hello Service}

\KeywordTok{[Service]}
\DataTypeTok{Type}\OtherTok{=}\StringTok{oneshot}
\DataTypeTok{ExecStart}\OtherTok{=}\StringTok{/usr/local/bin/hello.sh}

\KeywordTok{[Install]}
\DataTypeTok{WantedBy}\OtherTok{=}\StringTok{multi{-}user.target}
\end{Highlighting}
\end{Shaded}

Reload systemd (IMPORTANT anytime you create/edit/delete unit files):

\begin{Shaded}
\begin{Highlighting}[]
\FunctionTok{sudo}\NormalTok{ systemctl daemon{-}reload}
\FunctionTok{sudo}\NormalTok{ systemctl start hello.service}
\FunctionTok{cat}\NormalTok{ /tmp/hello.log    }\CommentTok{\# Boom \textbackslash{}faExclamationCircle}
\end{Highlighting}
\end{Shaded}

\begin{center}\rule{0.5\linewidth}{0.5pt}\end{center}

\paragraph{7) Editing \& Overriding Units (The Smart
Way)}\label{editing-overriding-units-the-smart-way}

\textbf{Never edit files in \texttt{/lib/systemd/system/}.}

Instead:

\begin{Shaded}
\begin{Highlighting}[]
\FunctionTok{sudo}\NormalTok{ systemctl edit apache2}
\end{Highlighting}
\end{Shaded}

This creates
\texttt{/etc/systemd/system/apache2.service.d/override.conf}. Safe
override.

\begin{Shaded}
\begin{Highlighting}[]
\KeywordTok{[Service]}
\DataTypeTok{Restart}\OtherTok{=}\StringTok{always}
\end{Highlighting}
\end{Shaded}

Save → \texttt{sudo\ systemctl\ daemon-reload} → restart service.

\begin{center}\rule{0.5\linewidth}{0.5pt}\end{center}

\paragraph{8) systemd Timers (Modern
Scheduling)}\label{systemd-timers-modern-scheduling}

Now we level up. \textbf{Timers = replacement for cron.}

\textbf{Why Use systemd Timers Instead of Cron?} - Integrated with
services - Better logging - More control - Dependency aware - Can use
calendar expressions

\begin{quote}
Note that systemd is more complex than Cron and requires systemd to work
on your system.
\end{quote}

\textbf{Timer Structure:} A timer requires (1) a \texttt{.service} file
and (2) a \texttt{.timer} file.

\begin{center}\rule{0.5\linewidth}{0.5pt}\end{center}

\paragraph{9) Example: Send Message to All
Users}\label{example-send-message-to-all-users}

Create script:

\begin{Shaded}
\begin{Highlighting}[]
\FunctionTok{sudo}\NormalTok{ nano /usr/local/bin/wall{-}msg.sh}
\end{Highlighting}
\end{Shaded}

\begin{Shaded}
\begin{Highlighting}[]
\CommentTok{\#!/bin/bash}
\FunctionTok{wall} \StringTok{"System reminder: Stay awesome!"}
\end{Highlighting}
\end{Shaded}

Make executable:

\begin{Shaded}
\begin{Highlighting}[]
\FunctionTok{sudo}\NormalTok{ chmod +x /usr/local/bin/wall{-}msg.sh}
\end{Highlighting}
\end{Shaded}

Create service file:

\begin{Shaded}
\begin{Highlighting}[]
\FunctionTok{sudo}\NormalTok{ nano /etc/systemd/system/wall{-}msg.service}
\end{Highlighting}
\end{Shaded}

\begin{Shaded}
\begin{Highlighting}[]
\KeywordTok{[Unit]}
\DataTypeTok{Description}\OtherTok{=}\StringTok{Wall Message Service}

\KeywordTok{[Service]}
\DataTypeTok{Type}\OtherTok{=}\StringTok{oneshot}
\DataTypeTok{ExecStart}\OtherTok{=}\StringTok{/usr/local/bin/wall{-}msg.sh}
\end{Highlighting}
\end{Shaded}

Create timer file:

\begin{Shaded}
\begin{Highlighting}[]
\FunctionTok{sudo}\NormalTok{ nano /etc/systemd/system/wall{-}msg.timer}
\end{Highlighting}
\end{Shaded}

\begin{Shaded}
\begin{Highlighting}[]
\KeywordTok{[Unit]}
\DataTypeTok{Description}\OtherTok{=}\StringTok{Run wall message every minute}

\KeywordTok{[Timer]}
\DataTypeTok{OnCalendar}\OtherTok{=}\StringTok{*{-}*{-}* *:*:00}
\DataTypeTok{Persistent}\OtherTok{=}\KeywordTok{true}

\KeywordTok{[Install]}
\DataTypeTok{WantedBy}\OtherTok{=}\StringTok{timers.target}
\end{Highlighting}
\end{Shaded}

\textbf{Breaking down the timer file:}

\textbf{\texttt{{[}Unit{]}} --- The ID Card} - \texttt{Description}: A
human-readable label. It doesn't affect logic, but when you run
\texttt{systemctl\ status}, this is what you see.

\textbf{\texttt{{[}Timer{]}} --- The Alarm Clock} - \texttt{OnCalendar}:
Defines \emph{when} the task runs (explained in detail below). -
\texttt{Persistent=true}: Acts like a ``snooze'' button for missed jobs.
If your computer was turned off when the timer was supposed to trigger,
systemd will \textbf{immediately run the task as soon as you turn the
computer back on}.

\textbf{\texttt{{[}Install{]}} --- The Connection} -
\texttt{WantedBy=timers.target}: Tells systemd how to ``hook'' this
timer into the system. When you run \texttt{systemctl\ enable}, it links
this timer to the standard group of system timers so it starts
automatically on boot.

Activate:

\begin{Shaded}
\begin{Highlighting}[]
\FunctionTok{sudo}\NormalTok{ systemctl daemon{-}reload}
\FunctionTok{sudo}\NormalTok{ systemctl enable wall{-}msg.timer}
\FunctionTok{sudo}\NormalTok{ systemctl start wall{-}msg.timer}
\ExtensionTok{systemctl}\NormalTok{ list{-}timers    }\CommentTok{\# Check timers \textbackslash{}faExclamationCircle}
\end{Highlighting}
\end{Shaded}

\begin{center}\rule{0.5\linewidth}{0.5pt}\end{center}

\paragraph{10) Understanding OnCalendar}\label{understanding-oncalendar}

\texttt{OnCalendar} uses a simple format:
\texttt{DayOfWeek\ Year-Month-Day\ Hour:Minute:Second}

\textbf{The Asterisk (\texttt{*}) means ``every'' or ``don't care''.}

{\def\LTcaptype{none} % do not increment counter
\begin{longtable}[]{@{}
  >{\raggedright\arraybackslash}p{(\linewidth - 4\tabcolsep) * \real{0.3333}}
  >{\raggedright\arraybackslash}p{(\linewidth - 4\tabcolsep) * \real{0.3333}}
  >{\raggedright\arraybackslash}p{(\linewidth - 4\tabcolsep) * \real{0.3333}}@{}}
\toprule\noalign{}
\begin{minipage}[b]{\linewidth}\raggedright
Expression
\end{minipage} & \begin{minipage}[b]{\linewidth}\raggedright
Meaning
\end{minipage} & \begin{minipage}[b]{\linewidth}\raggedright
Breakdown
\end{minipage} \\
\midrule\noalign{}
\endhead
\bottomrule\noalign{}
\endlastfoot
\texttt{*-*-*\ *:*:00} & Every minute & Every year-month-day, every
hour:minute, at second 00 \\
\texttt{daily} & Once per day & Shortcut for \texttt{*-*-*\ 00:00:00} \\
\texttt{weekly} & Once per week & Shortcut for
\texttt{Mon\ *-*-*\ 00:00:00} \\
\texttt{Mon\ *-*-*\ 09:00:00} & Every Monday at 9AM & Monday, any date,
at 09:00:00 \\
\texttt{*-*-01\ 00:00:00} & First day of every month & Any year-month,
day 01, at midnight \\
\texttt{*-*-*\ *:00:00} & Every hour, on the hour & Any date, every
hour, at minute 00, second 00 \\
\texttt{Mon..Fri\ *-*-*\ 08:30:00} & Weekdays at 8:30AM & Monday through
Friday, any date, at 08:30 \\
\texttt{*-*-12,24\ 03:00:00} & 12th and 24th of every month at 3AM & Any
year-month, days 12 and 24, at 03:00 \\
\end{longtable}
}

\textbf{Pro Tips:} - You can use shortcuts: \texttt{hourly},
\texttt{daily}, \texttt{weekly}, \texttt{monthly}, \texttt{yearly} - Use
\texttt{systemd-analyze\ calendar\ "YOUR\_EXPRESSION"} to test and
verify your schedule before using it.

\begin{Shaded}
\begin{Highlighting}[]
\CommentTok{\# Test a calendar expression}
\ExtensionTok{systemd{-}analyze}\NormalTok{ calendar }\StringTok{"Mon *{-}*{-}* 09:00:00"}
\end{Highlighting}
\end{Shaded}

This command will show you the next scheduled run time so you can verify
your expression is correct.

\begin{center}\rule{0.5\linewidth}{0.5pt}\end{center}

\paragraph{11) Other Timer Options}\label{other-timer-options}

\textbf{Other Timer Options:} - \texttt{OnBootSec=} --- After boot -
\texttt{OnUnitActiveSec=} --- After last run - \texttt{Persistent=true}
--- Run missed jobs after reboot

\begin{center}\rule{0.5\linewidth}{0.5pt}\end{center}

\paragraph{12) Timer vs Cron Comparison}\label{timer-vs-cron-comparison}

{\def\LTcaptype{none} % do not increment counter
\begin{longtable}[]{@{}lll@{}}
\toprule\noalign{}
Feature & Cron & systemd Timer \\
\midrule\noalign{}
\endhead
\bottomrule\noalign{}
\endlastfoot
Logging & Basic & journalctl \\
Dependencies & No & Yes \\
Missed Runs & No & Yes (Persistent) \\
Integration & Separate & Native \\
\end{longtable}
}

\textbf{Cron jobs for repetitive tasks:}

\begin{Shaded}
\begin{Highlighting}[]
\FunctionTok{crontab} \AttributeTok{{-}l}      \CommentTok{\# List cron jobs}
\FunctionTok{crontab} \AttributeTok{{-}e}      \CommentTok{\# Edit cron jobs}
\end{Highlighting}
\end{Shaded}

\begin{center}\rule{0.5\linewidth}{0.5pt}\end{center}

\paragraph{13) Debugging Services \&
Timers}\label{debugging-services-timers}

Check service logs and failures:

\begin{Shaded}
\begin{Highlighting}[]
\ExtensionTok{journalctl} \AttributeTok{{-}u}\NormalTok{ wall{-}msg.service    }\CommentTok{\# Logs for specific service}
\ExtensionTok{journalctl} \AttributeTok{{-}f}                     \CommentTok{\# Live logs}
\ExtensionTok{systemctl} \AttributeTok{{-}{-}failed}                \CommentTok{\# Check failures}
\end{Highlighting}
\end{Shaded}

\begin{center}\rule{0.5\linewidth}{0.5pt}\end{center}

\paragraph{14) Real World Use Cases}\label{real-world-use-cases}

\begin{itemize}
\tightlist
\item
  Backup scripts
\item
  Log cleanup
\item
  Health checks
\item
  Auto-updates
\item
  Dev server restarts
\item
  IoT tasks
\end{itemize}

\textbf{Common Mistakes:} - \faTimesCircle Forgetting
\texttt{daemon-reload} - \faTimesCircle Editing
\texttt{/lib/systemd/system/} directly - \faTimesCircle Wrong file
permissions - \faTimesCircle Forgetting \texttt{WantedBy=}

\begin{center}\rule{0.5\linewidth}{0.5pt}\end{center}

\textbf{Final Mental Model:} - Service = What to run - Timer = When to
run - Unit file = Configuration blueprint - systemctl = Control center

\begin{quote}
You are no longer a Linux user. You're becoming a Linux operator. →
\end{quote}

\begin{center}\rule{0.5\linewidth}{0.5pt}\end{center}

\paragraph{Useful Links --- Session 8}\label{useful-links-session-8}

\textbf{systemd Fundamentals:}

%
  \par\vspace{0.5ex}%
  \href{https://www.geeksforgeeks.org/linux-unix/linux-systemd-and-its-components/}{\textcolor{blue}{\underline{\faLink\quad Linux systemd and its components (GeeksforGeeks)}}}%
  \par\vspace{0.5ex}%

%
  \par\vspace{0.5ex}%
  \href{https://www.geeksforgeeks.org/linux-unix/an-introduction-to-systemd-and-its-role-in-the-boot-process/}{\textcolor{blue}{\underline{\faLink\quad Introduction to systemd (GeeksforGeeks)}}}%
  \par\vspace{0.5ex}%

%
  \par\vspace{0.5ex}%
  \href{https://www.geeksforgeeks.org/linux-unix/how-to-mask-a-systemd-unit-in-linux/}{\textcolor{blue}{\underline{\faLink\quad How to mask a systemd unit (GeeksforGeeks)}}}%
  \par\vspace{0.5ex}%

\begin{center}\rule{0.5\linewidth}{0.5pt}\end{center}

\subsubsection{Session 9: Storage Management \& Performance
Optimization}\label{session-9-storage-management-performance-optimization}

\begin{quote}
\textbf{\texttt{zram-tools} is optional but recommended for low-memory
systems.}
\end{quote}

\paragraph{Learning outcomes}\label{learning-outcomes-8}

\begin{itemize}
\tightlist
\item
  Inspect and mount storage devices correctly
\item
  Configure persistent mounts through \texttt{/etc/fstab}
\item
  Monitor memory behavior and choose swap/zram strategy
\item
  Manage datasets/models under embedded storage constraints
\end{itemize}

\begin{center}\rule{0.5\linewidth}{0.5pt}\end{center}

\paragraph{1) Why This Session Matters}\label{why-this-session-matters}

Before we touch commands, understand this: - Your \textbf{disk} stores
data permanently - Your \textbf{RAM} is fast temporary memory -
\textbf{Swap} is emergency backup memory - \textbf{zram} is compressed
RAM swap (smart trick for low-memory systems)

If you understand this session, you'll mount extra drives safely, fix
broken boot issues, add a second disk like a pro, understand RAM
pressure, and boost low-RAM machines using zram.

\textbf{This is not theory. This is real Linux power.}

\begin{center}\rule{0.5\linewidth}{0.5pt}\end{center}

\paragraph{2) Understanding Your Disks}\label{understanding-your-disks}

\textbf{\texttt{lsblk} --- The Disk Tree Viewer:}

\begin{Shaded}
\begin{Highlighting}[]
\ExtensionTok{lsblk}
\ExtensionTok{lsblk} \AttributeTok{{-}f}          \CommentTok{\# Shows filesystem, UUID, mountpoint (MORE DETAILED)}
\end{Highlighting}
\end{Shaded}

Example output:

\begin{verbatim}
sda
├─sda1
├─sda2
└─sda3
sdb
└─sdb1
\end{verbatim}

\textbf{Shows:} Disk names, partitions, mount points, size, type.\\
\faExclamationCircle Use this before and after plugging a USB.

\textbf{\texttt{blkid} --- UUID Finder:}

\begin{Shaded}
\begin{Highlighting}[]
\FunctionTok{sudo}\NormalTok{ blkid}
\end{Highlighting}
\end{Shaded}

\begin{verbatim}
/dev/sda1: UUID="1A2B-3C4D" TYPE="vfat"
\end{verbatim}

Why important? Because in \texttt{/etc/fstab}, we use \textbf{UUID}, not
\texttt{/dev/sda1}. Why? Because device names can change on boot.
\textbf{UUID never lies.}

\begin{center}\rule{0.5\linewidth}{0.5pt}\end{center}

\paragraph{3) Mounting and Unmounting
Drives}\label{mounting-and-unmounting-drives}

Unlike Windows, Linux does not auto-use disks unless mounted.

\textbf{Manual Mount:}

\begin{Shaded}
\begin{Highlighting}[]
\FunctionTok{sudo}\NormalTok{ mkdir /mnt/mydrive           }\CommentTok{\# Step 1: Create mount point}
\FunctionTok{sudo}\NormalTok{ mount /dev/sdb1 /mnt/mydrive }\CommentTok{\# Step 2: Mount}
\FunctionTok{ls}\NormalTok{ /mnt/mydrive                   }\CommentTok{\# Step 3: Check}
\end{Highlighting}
\end{Shaded}

\textbf{Unmount:}

\begin{Shaded}
\begin{Highlighting}[]
\FunctionTok{sudo}\NormalTok{ umount /mnt/mydrive}
\CommentTok{\# OR}
\FunctionTok{sudo}\NormalTok{ umount /dev/sdb1}
\end{Highlighting}
\end{Shaded}

\faExclamationCircle If busy error:

\begin{Shaded}
\begin{Highlighting}[]
\ExtensionTok{lsof} \KeywordTok{|} \FunctionTok{grep}\NormalTok{ sdb1    }\CommentTok{\# Find what\textquotesingle{}s using it}
\end{Highlighting}
\end{Shaded}

\begin{center}\rule{0.5\linewidth}{0.5pt}\end{center}

\paragraph{\texorpdfstring{4) Persistent Mounts with
\texttt{/etc/fstab}}{4) Persistent Mounts with /etc/fstab}}\label{persistent-mounts-with-etcfstab}

This file controls what mounts at boot.

\begin{Shaded}
\begin{Highlighting}[]
\FunctionTok{sudo}\NormalTok{ nano /etc/fstab}
\end{Highlighting}
\end{Shaded}

Example entry:

\begin{verbatim}
UUID=1a2b3c4d-xxxx-xxxx-xxxx-xxxxxxxxxxxx  /mnt/mydrive  ext4  defaults  0  2
\end{verbatim}

{\def\LTcaptype{none} % do not increment counter
\begin{longtable}[]{@{}ll@{}}
\toprule\noalign{}
Column & Meaning \\
\midrule\noalign{}
\endhead
\bottomrule\noalign{}
\endlastfoot
UUID & Disk identifier \\
Mount point & Where it appears \\
Filesystem & ext4, ntfs, vfat \\
Options & defaults, noatime \\
Dump & usually 0 \\
fsck order & usually 2 \\
\end{longtable}
}

\textbf{{[}LAB{]} Test Before Reboot:}

\begin{Shaded}
\begin{Highlighting}[]
\FunctionTok{sudo}\NormalTok{ mount }\AttributeTok{{-}a}    \CommentTok{\# If no errors → safe to reboot. If error → FIX IT before rebooting.}
\end{Highlighting}
\end{Shaded}

\begin{quote}
\textbf{Pro tip:} Use \texttt{nofail} option in fstab for external
drives so system doesn't panic if drive is missing.
\end{quote}

\begin{center}\rule{0.5\linewidth}{0.5pt}\end{center}

\paragraph{5) Monitoring RAM Usage}\label{monitoring-ram-usage}

RAM = short-term memory.

\begin{Shaded}
\begin{Highlighting}[]
\FunctionTok{free} \AttributeTok{{-}h}
\end{Highlighting}
\end{Shaded}

Example:

\begin{verbatim}
              total   used   free  shared  buff/cache  available
Mem:           7.6G   3.1G   1.2G     500M     3.3G       3.8G
Swap:          2.0G   0B     2.0G
\end{verbatim}

\textbf{Important columns:} - \textbf{used} → currently used -
\textbf{available} → real usable memory - \textbf{buff/cache} → Linux
using RAM smartly

\faExclamationCircle Linux uses RAM aggressively. That's GOOD. Unused
RAM = wasted RAM.

\textbf{Real-time view:}

\begin{Shaded}
\begin{Highlighting}[]
\ExtensionTok{watch} \AttributeTok{{-}n}\NormalTok{ 1 free }\AttributeTok{{-}h}    \CommentTok{\# Updates every second}
\end{Highlighting}
\end{Shaded}

\begin{center}\rule{0.5\linewidth}{0.5pt}\end{center}

\paragraph{6) Understanding Swap}\label{understanding-swap}

Swap = disk used as overflow memory. When RAM fills up → inactive pages
move to swap. But swap is slower than RAM.

\begin{Shaded}
\begin{Highlighting}[]
\ExtensionTok{swapon} \AttributeTok{{-}{-}show}         \CommentTok{\# Check swap}
\FunctionTok{cat}\NormalTok{ /proc/swaps}
\end{Highlighting}
\end{Shaded}

\textbf{Creating a Swap File (Manual Way):}

\begin{Shaded}
\begin{Highlighting}[]
\FunctionTok{sudo}\NormalTok{ fallocate }\AttributeTok{{-}l}\NormalTok{ 2G /swapfile    }\CommentTok{\# Step 1}
\FunctionTok{sudo}\NormalTok{ chmod 600 /swapfile          }\CommentTok{\# Step 2}
\FunctionTok{sudo}\NormalTok{ mkswap /swapfile             }\CommentTok{\# Step 3}
\FunctionTok{sudo}\NormalTok{ swapon /swapfile             }\CommentTok{\# Step 4}
\end{Highlighting}
\end{Shaded}

Make permanent---add to \texttt{/etc/fstab}:

\begin{verbatim}
/swapfile none swap sw 0 0
\end{verbatim}

\begin{center}\rule{0.5\linewidth}{0.5pt}\end{center}

\paragraph{7) Swap and zram}\label{swap-and-zram}

\textbf{zram --- Compressed RAM Swap:}\\
zram creates compressed swap \emph{in RAM}. Instead of writing to disk,
it compresses pages and stores them in RAM---faster than disk swap.

Great for 4GB RAM laptops, old machines, and low-resource VPS.

\textbf{Install on Debian/Ubuntu/Mint:}

\begin{Shaded}
\begin{Highlighting}[]
\FunctionTok{sudo}\NormalTok{ apt install zram{-}tools}
\ExtensionTok{dpkg} \AttributeTok{{-}l} \KeywordTok{|} \FunctionTok{grep}\NormalTok{ zram              }\CommentTok{\# Check if installed}
\ExtensionTok{systemctl}\NormalTok{ status zramswap        }\CommentTok{\# Check if active}
\end{Highlighting}
\end{Shaded}

Check config:

\begin{Shaded}
\begin{Highlighting}[]
\FunctionTok{cat}\NormalTok{ /etc/default/zramswap}
\end{Highlighting}
\end{Shaded}

You can adjust \texttt{PERCENT=70} (use 70\% of RAM for
zram---recommended: 70\%-80\%).

\begin{Shaded}
\begin{Highlighting}[]
\FunctionTok{sudo}\NormalTok{ systemctl restart zramswap}
\ExtensionTok{swapon} \AttributeTok{{-}{-}show}                    \CommentTok{\# Should show /dev/zram0}
\end{Highlighting}
\end{Shaded}

\begin{center}\rule{0.5\linewidth}{0.5pt}\end{center}

\textbf{Swap vs zram Comparison:}

{\def\LTcaptype{none} % do not increment counter
\begin{longtable}[]{@{}lll@{}}
\toprule\noalign{}
Feature & Disk Swap & zram \\
\midrule\noalign{}
\endhead
\bottomrule\noalign{}
\endlastfoot
Speed & Slow & Faster \\
Uses disk & Yes & No \\
Best for & Large systems & Low RAM \\
Compression & No & Yes \\
\end{longtable}
}

\begin{quote}
\textbf{Best practice for laptops:} Use zram + keep small disk swap.
\end{quote}

\begin{center}\rule{0.5\linewidth}{0.5pt}\end{center}

\paragraph{8) Embedded Storage Choices (SD vs
SSD)}\label{embedded-storage-choices-sd-vs-ssd}

\begin{itemize}
\tightlist
\item
  \textbf{SD:} Portable and common, but limited endurance
\item
  \textbf{SSD:} Better durability/performance for heavy read/write
  workloads
\end{itemize}

\begin{center}\rule{0.5\linewidth}{0.5pt}\end{center}

\paragraph{9) Storage Lifespan
Optimization}\label{storage-lifespan-optimization}

Reduce unnecessary writes, rotate logs, and separate high-write data
paths where possible.

\begin{center}\rule{0.5\linewidth}{0.5pt}\end{center}

\paragraph{10) AI and Simulation Storage
Management}\label{ai-and-simulation-storage-management}

Organize datasets/models/simulation outputs with clear cleanup and
archival policy.

\begin{center}\rule{0.5\linewidth}{0.5pt}\end{center}

\paragraph{11) Troubleshooting Boot Failures (fstab
Mistakes)}\label{troubleshooting-boot-failures-fstab-mistakes}

If system won't boot: 1. Boot into recovery mode 2. Mount root 3. Edit
\texttt{/etc/fstab} 4. Fix or comment broken line

\begin{center}\rule{0.5\linewidth}{0.5pt}\end{center}

\paragraph{12) Useful Extra Commands}\label{useful-extra-commands}

Additional disk and memory utilities:

\begin{Shaded}
\begin{Highlighting}[]
\FunctionTok{df} \AttributeTok{{-}h}                 \CommentTok{\# Disk usage}
\FunctionTok{du} \AttributeTok{{-}sh} \PreprocessorTok{*}              \CommentTok{\# Folder sizes}
\FunctionTok{mount} \KeywordTok{|} \FunctionTok{grep}\NormalTok{ sdb      }\CommentTok{\# See mount points}
\end{Highlighting}
\end{Shaded}

\begin{center}\rule{0.5\linewidth}{0.5pt}\end{center}

\paragraph{Practice}\label{practice-1}

Verify your disk and memory configuration:

\begin{Shaded}
\begin{Highlighting}[]
\ExtensionTok{lsblk}
\ExtensionTok{blkid}
\FunctionTok{free} \AttributeTok{{-}h}
\FunctionTok{vmstat}\NormalTok{ 1 5}
\FunctionTok{df} \AttributeTok{{-}h}
\end{Highlighting}
\end{Shaded}

\begin{center}\rule{0.5\linewidth}{0.5pt}\end{center}

\paragraph{Practical Lab Exercises}\label{practical-lab-exercises}

\begin{enumerate}
\def\labelenumi{\arabic{enumi}.}
\tightlist
\item
  Plug USB → identify it using \texttt{lsblk}
\item
  Mount it manually in \texttt{/mnt/testusb}
\item
  Make it permanent in fstab
\item
  Create 1GB swap file
\item
  Install zram and compare swap usage
\end{enumerate}

\begin{quote}
\textbf{Final Advice:} Disk and memory management is not optional
knowledge. It's what separates ``Linux user'' from ``Linux operator''.
Break things in a VM. Test everything with \texttt{mount\ -a}. Never
reboot blindly after editing fstab. You're building serious Linux skills
now.
\end{quote}

\begin{center}\rule{0.5\linewidth}{0.5pt}\end{center}

\paragraph{Useful Links --- Session 9}\label{useful-links-session-9}

\textbf{Disk \& Storage Tutorials:}

%
  \par\vspace{0.5ex}%
  \href{https://mmbesar.github.io/Tutorials/2nd-drive/}{\textcolor{blue}{\underline{\faLink\quad Adding 2nd drive tutorial (mmbesar)}}}%
  \par\vspace{0.5ex}%

%
  \par\vspace{0.5ex}%
  \href{https://mmbesar.github.io/Tutorials/RAM-SWAP-ZRAM/}{\textcolor{blue}{\underline{\faLink\quad RAM, SWAP, ZRAM tutorial (mmbesar)}}}%
  \par\vspace{0.5ex}%

%
  \par\vspace{0.5ex}%
  \href{https://www.youtube.com/watch?v=yWuHI7uoftY}{\textcolor{blue}{\underline{\faLink\quad Disk management video guide}}}%
  \par\vspace{0.5ex}%

%
  \par\vspace{0.5ex}%
  \href{https://youtu.be/CPvJvY79vEw}{\textcolor{blue}{\underline{\faLink\quad Swap \& ZRAM video guide}}}%
  \par\vspace{0.5ex}%

\begin{center}\rule{0.5\linewidth}{0.5pt}\end{center}

\subsubsection{Session 10: Shell Scripting \& Robotics
Networking}\label{session-10-shell-scripting-robotics-networking}

\begin{quote}
\textbf{Verify packages:} \texttt{curl}, \texttt{aria2}, \texttt{ssh},
\texttt{net-tools}
\end{quote}

\paragraph{Learning outcomes}\label{learning-outcomes-9}

\begin{itemize}
\tightlist
\item
  Build practical Bash scripts for robotics operations
\item
  Diagnose and configure network connectivity
\item
  Transfer data/models securely and efficiently
\item
  Set up static IP and SSH keys for automation
\end{itemize}

\begin{center}\rule{0.5\linewidth}{0.5pt}\end{center}

\subsubsection{Part 1: Bash Scripting for Robotics
Automation}\label{part-1-bash-scripting-for-robotics-automation}

You already know how to navigate the terminal. You've used
\texttt{grep}, \texttt{sed}, and \texttt{awk}. You've redirected output
and written basic scripts. Now it's time to think like a developer and
build tools that actually \emph{do} things.

\begin{center}\rule{0.5\linewidth}{0.5pt}\end{center}

\paragraph{1) Variables and Data Types}\label{variables-and-data-types}

In Bash, everything is a string. But that doesn't mean you can't be
clever.

\begin{Shaded}
\begin{Highlighting}[]
\CommentTok{\#!/bin/bash}

\CommentTok{\# Variable assignment (no spaces!)}
\VariableTok{user\_name}\OperatorTok{=}\StringTok{"hacker"}
\VariableTok{age}\OperatorTok{=}\NormalTok{21}
\VariableTok{is\_cool}\OperatorTok{=}\NormalTok{true}

\CommentTok{\# Accessing variables}
\BuiltInTok{echo} \StringTok{"Hey }\VariableTok{$user\_name}\StringTok{, you\textquotesingle{}re }\VariableTok{$age}\StringTok{ years old!"}

\CommentTok{\# Command substitution}
\VariableTok{current\_date}\OperatorTok{=}\VariableTok{$(}\FunctionTok{date}\NormalTok{ +\%Y{-}\%m{-}\%d}\VariableTok{)}
\VariableTok{kernel\_version}\OperatorTok{=}\VariableTok{$(}\FunctionTok{uname} \AttributeTok{{-}r}\VariableTok{)}

\BuiltInTok{echo} \StringTok{"Today is }\VariableTok{$current\_date}\StringTok{, running kernel }\VariableTok{$kernel\_version}\StringTok{"}
\end{Highlighting}
\end{Shaded}

\begin{quote}
\textbf{Pro Tip:} Use \texttt{\$()} for command substitution instead of
backticks. It's nestable and easier to read.
\end{quote}

\textbf{Core script blocks to practice:} Variables, Conditions and
loops, Functions, Exit codes and basic error handling.

\begin{center}\rule{0.5\linewidth}{0.5pt}\end{center}

\paragraph{2) Control Flow}\label{control-flow}

\textbf{If Statements:}

\begin{Shaded}
\begin{Highlighting}[]
\CommentTok{\#!/bin/bash}

\VariableTok{file\_count}\OperatorTok{=}\VariableTok{$(}\FunctionTok{ls} \KeywordTok{|} \FunctionTok{wc} \AttributeTok{{-}l}\VariableTok{)}

\ControlFlowTok{if} \BuiltInTok{[} \VariableTok{$file\_count} \OtherTok{{-}gt}\NormalTok{ 10 }\BuiltInTok{]}\KeywordTok{;} \ControlFlowTok{then}
    \BuiltInTok{echo} \StringTok{"Directory is cluttered (}\VariableTok{$file\_count}\StringTok{ files). Time to clean up!"}
\ControlFlowTok{elif} \BuiltInTok{[} \VariableTok{$file\_count} \OtherTok{{-}eq}\NormalTok{ 0 }\BuiltInTok{]}\KeywordTok{;} \ControlFlowTok{then}
    \BuiltInTok{echo} \StringTok{"Empty directory. Lonely vibes."}
\ControlFlowTok{else}
    \BuiltInTok{echo} \StringTok{"Looking good with }\VariableTok{$file\_count}\StringTok{ files."}
\ControlFlowTok{fi}
\end{Highlighting}
\end{Shaded}

\begin{quote}
\textbf{Important:} Always use spaces inside brackets \texttt{{[}\ {]}}.
\texttt{{[}\ \$var\ -eq\ 1\ {]}} works, \texttt{{[}\$var\ -eq\ 1{]}}
fails.
\end{quote}

\textbf{Case Statements:}

\begin{Shaded}
\begin{Highlighting}[]
\CommentTok{\#!/bin/bash}

\BuiltInTok{echo} \StringTok{"Choose your distro:"}
\BuiltInTok{echo} \StringTok{"1) Arch"}
\BuiltInTok{echo} \StringTok{"2) Debian"}
\BuiltInTok{echo} \StringTok{"3) Fedora"}
\BuiltInTok{read} \VariableTok{choice}

\ControlFlowTok{case} \VariableTok{$choice} \KeywordTok{in}
    \SpecialStringTok{1}\KeywordTok{)} \BuiltInTok{echo} \StringTok{"I use Arch, btw"} \ControlFlowTok{;;}
    \SpecialStringTok{2}\KeywordTok{)} \BuiltInTok{echo} \StringTok{"Stable and reliable. Nice."} \ControlFlowTok{;;}
    \SpecialStringTok{3}\KeywordTok{)} \BuiltInTok{echo} \StringTok{"Red Hat family represent!"} \ControlFlowTok{;;}
    \PreprocessorTok{*}\KeywordTok{)} \BuiltInTok{echo} \StringTok{"Invalid choice. Try again."} \ControlFlowTok{;;}
\ControlFlowTok{esac}
\end{Highlighting}
\end{Shaded}

\begin{center}\rule{0.5\linewidth}{0.5pt}\end{center}

\paragraph{3) Loops}\label{loops}

\textbf{For Loops:}

\begin{Shaded}
\begin{Highlighting}[]
\CommentTok{\#!/bin/bash}

\ControlFlowTok{for}\NormalTok{ file }\KeywordTok{in} \PreprocessorTok{*}\NormalTok{.txt}\KeywordTok{;} \ControlFlowTok{do}
    \BuiltInTok{echo} \StringTok{"Processing }\VariableTok{$file}\StringTok{..."}
    \FunctionTok{wc} \AttributeTok{{-}l} \StringTok{"}\VariableTok{$file}\StringTok{"}
\ControlFlowTok{done}

\CommentTok{\# C{-}style}
\ControlFlowTok{for} \KeywordTok{((}\VariableTok{i}\OperatorTok{=}\DecValTok{0}\KeywordTok{;} \VariableTok{i}\OperatorTok{\textless{}}\DecValTok{5}\KeywordTok{;} \VariableTok{i}\OperatorTok{++}\KeywordTok{));} \ControlFlowTok{do}
    \BuiltInTok{echo} \StringTok{"Iteration }\VariableTok{$i}\StringTok{"}
\ControlFlowTok{done}
\end{Highlighting}
\end{Shaded}

\textbf{While Loops:}

\begin{Shaded}
\begin{Highlighting}[]
\CommentTok{\#!/bin/bash}

\CommentTok{\# Read file line by line}
\ControlFlowTok{while} \VariableTok{IFS}\OperatorTok{=} \BuiltInTok{read} \AttributeTok{{-}r} \VariableTok{line}\KeywordTok{;} \ControlFlowTok{do}
    \BuiltInTok{echo} \StringTok{"Line: }\VariableTok{$line}\StringTok{"}
\ControlFlowTok{done} \OperatorTok{\textless{}}\NormalTok{ input.txt}

\CommentTok{\# With break}
\VariableTok{counter}\OperatorTok{=}\NormalTok{0}
\ControlFlowTok{while} \FunctionTok{true}\KeywordTok{;} \ControlFlowTok{do}
    \BuiltInTok{echo} \StringTok{"Loop }\VariableTok{$counter}\StringTok{"}
    \KeywordTok{((}\VariableTok{counter}\OperatorTok{++}\KeywordTok{))}
    \BuiltInTok{[} \VariableTok{$counter} \OtherTok{{-}eq}\NormalTok{ 5 }\BuiltInTok{]} \KeywordTok{\&\&} \ControlFlowTok{break}
\ControlFlowTok{done}
\end{Highlighting}
\end{Shaded}

\begin{center}\rule{0.5\linewidth}{0.5pt}\end{center}

\paragraph{4) Functions}\label{functions}

Write reusable code blocks:

\begin{Shaded}
\begin{Highlighting}[]
\CommentTok{\#!/bin/bash}

\FunctionTok{check\_port()} \KeywordTok{\{}
    \BuiltInTok{local} \VariableTok{host}\OperatorTok{=}\VariableTok{$1}
    \BuiltInTok{local} \VariableTok{port}\OperatorTok{=}\VariableTok{$2}
    \FunctionTok{timeout}\NormalTok{ 2 bash }\AttributeTok{{-}c} \StringTok{"\textless{}/dev/tcp/}\VariableTok{$host}\StringTok{/}\VariableTok{$port}\StringTok{"} \DecValTok{2}\OperatorTok{\textgreater{}}\NormalTok{/dev/null }\KeywordTok{\&\&} \DataTypeTok{\textbackslash{}}
        \BuiltInTok{echo} \StringTok{"Port }\VariableTok{$port}\StringTok{ on }\VariableTok{$host}\StringTok{ is OPEN"} \KeywordTok{||} \DataTypeTok{\textbackslash{}}
        \BuiltInTok{echo} \StringTok{"Port }\VariableTok{$port}\StringTok{ on }\VariableTok{$host}\StringTok{ is CLOSED"}
\KeywordTok{\}}

\ExtensionTok{check\_port} \StringTok{"google.com"}\NormalTok{ 443}
\ExtensionTok{check\_port} \StringTok{"localhost"}\NormalTok{ 22}
\end{Highlighting}
\end{Shaded}

\textbf{Key Points:} - Use \texttt{local} for variables inside functions
to avoid global scope pollution - Functions should do one thing well
(Unix philosophy) - Return values are 0-255 (exit codes), not strings

\begin{center}\rule{0.5\linewidth}{0.5pt}\end{center}

\paragraph{5) Error Handling}\label{error-handling}

Don't let your scripts crash silently:

\begin{Shaded}
\begin{Highlighting}[]
\CommentTok{\#!/bin/bash}

\CommentTok{\# Automatic error handling}
\BuiltInTok{set} \AttributeTok{{-}euo}\NormalTok{ pipefail}
\CommentTok{\# {-}e: Exit on error}
\CommentTok{\# {-}u: Exit on undefined variable}
\CommentTok{\# {-}o pipefail: Catch errors in pipelines}

\CommentTok{\# Or handle manually}
\ControlFlowTok{if} \OtherTok{! }\FunctionTok{ping} \AttributeTok{{-}c}\NormalTok{ 1 google.com }\OperatorTok{\&\textgreater{}}\NormalTok{ /dev/null}\KeywordTok{;} \ControlFlowTok{then}
    \BuiltInTok{echo} \StringTok{"ERROR: No internet connection"} \OperatorTok{\textgreater{}\&}\DecValTok{2}
    \BuiltInTok{exit}\NormalTok{ 1}
\ControlFlowTok{fi}
\end{Highlighting}
\end{Shaded}

\begin{center}\rule{0.5\linewidth}{0.5pt}\end{center}

\paragraph{6) Working with APIs}\label{working-with-apis}

Modern scripting is all about data:

\begin{Shaded}
\begin{Highlighting}[]
\CommentTok{\#!/bin/bash}

\VariableTok{response}\OperatorTok{=}\VariableTok{$(}\ExtensionTok{curl} \AttributeTok{{-}s} \StringTok{"https://api.github.com/users/torvalds"}\VariableTok{)}
\VariableTok{login}\OperatorTok{=}\VariableTok{$(}\BuiltInTok{echo} \StringTok{"}\VariableTok{$response}\StringTok{"} \KeywordTok{|} \ExtensionTok{jq} \AttributeTok{{-}r} \StringTok{\textquotesingle{}.login\textquotesingle{}}\VariableTok{)}
\VariableTok{followers}\OperatorTok{=}\VariableTok{$(}\BuiltInTok{echo} \StringTok{"}\VariableTok{$response}\StringTok{"} \KeywordTok{|} \ExtensionTok{jq} \AttributeTok{{-}r} \StringTok{\textquotesingle{}.followers\textquotesingle{}}\VariableTok{)}
\VariableTok{created\_at}\OperatorTok{=}\VariableTok{$(}\BuiltInTok{echo} \StringTok{"}\VariableTok{$response}\StringTok{"} \KeywordTok{|} \ExtensionTok{jq} \AttributeTok{{-}r} \StringTok{\textquotesingle{}.created\_at\textquotesingle{}}\VariableTok{)}

\BuiltInTok{echo} \StringTok{"User: }\VariableTok{$login}\StringTok{"}
\BuiltInTok{echo} \StringTok{"Followers: }\VariableTok{$followers}\StringTok{"}
\BuiltInTok{echo} \StringTok{"Joined: }\VariableTok{$created\_at}\StringTok{"}
\end{Highlighting}
\end{Shaded}

\begin{center}\rule{0.5\linewidth}{0.5pt}\end{center}

\subsubsection{Part 2: Networking
Essentials}\label{part-2-networking-essentials}

Linux networking tools are your eyes and ears into the digital world.
Whether you're debugging connectivity, downloading files, or managing
remote systems, these commands are essential.

\begin{center}\rule{0.5\linewidth}{0.5pt}\end{center}

\paragraph{1) Connectivity Testing}\label{connectivity-testing}

\textbf{ping --- The Classic:}

\begin{Shaded}
\begin{Highlighting}[]
\FunctionTok{ping}\NormalTok{ google.com}
\FunctionTok{ping} \AttributeTok{{-}c}\NormalTok{ 4 8.8.8.8           }\CommentTok{\# Stop after 4 packets}
\FunctionTok{sudo}\NormalTok{ ping }\AttributeTok{{-}i}\NormalTok{ 0.2 google.com }\CommentTok{\# Faster ping}
\FunctionTok{ping} \AttributeTok{{-}a}\NormalTok{ 192.168.1.1         }\CommentTok{\# Audible ping}
\end{Highlighting}
\end{Shaded}

\textbf{traceroute / tracepath:}

\begin{Shaded}
\begin{Highlighting}[]
\ExtensionTok{traceroute}\NormalTok{ google.com       }\CommentTok{\# See the path your packets take}
\ExtensionTok{tracepath}\NormalTok{ google.com        }\CommentTok{\# Alternative without root}
\end{Highlighting}
\end{Shaded}

\begin{center}\rule{0.5\linewidth}{0.5pt}\end{center}

\paragraph{2) Network Interface
Configuration}\label{network-interface-configuration}

\textbf{ip (Modern Replacement for ifconfig):}

\begin{Shaded}
\begin{Highlighting}[]
\ExtensionTok{ip}\NormalTok{ a                        }\CommentTok{\# Show all interfaces (short form)}
\ExtensionTok{ip}\NormalTok{ addr show eth0           }\CommentTok{\# Specific interface}
\FunctionTok{sudo}\NormalTok{ ip link set eth0 up    }\CommentTok{\# Bring interface up/down}
\FunctionTok{sudo}\NormalTok{ ip link set eth0 down}
\FunctionTok{sudo}\NormalTok{ ip addr add 192.168.1.100/24 dev eth0  }\CommentTok{\# Add IP}
\ExtensionTok{ip}\NormalTok{ route show               }\CommentTok{\# Routing table}
\FunctionTok{sudo}\NormalTok{ ip route add default via 192.168.1.1   }\CommentTok{\# Add route}
\end{Highlighting}
\end{Shaded}

\textbf{Why \texttt{ip} over \texttt{ifconfig}?} \texttt{ip} is part of
the iproute2 package, actively maintained, and supports modern
networking features.

\textbf{ifconfig (Legacy---still works on many systems):}

\begin{Shaded}
\begin{Highlighting}[]
\ExtensionTok{ifconfig}
\ExtensionTok{ifconfig}\NormalTok{ eth0}
\end{Highlighting}
\end{Shaded}

\begin{center}\rule{0.5\linewidth}{0.5pt}\end{center}

\paragraph{3) Data Transfer}\label{data-transfer}

\textbf{curl --- The URL Swiss Army Knife:}

\begin{Shaded}
\begin{Highlighting}[]
\ExtensionTok{curl}\NormalTok{ https://api.example.com/data                }\CommentTok{\# GET request}
\ExtensionTok{curl} \AttributeTok{{-}o}\NormalTok{ file.zip https://example.com/file.zip    }\CommentTok{\# Save to file}
\ExtensionTok{curl} \AttributeTok{{-}L}\NormalTok{ https://bit.ly/xyz                       }\CommentTok{\# Follow redirects}
\ExtensionTok{curl} \AttributeTok{{-}I}\NormalTok{ https://google.com                       }\CommentTok{\# Headers only}
\ExtensionTok{curl} \AttributeTok{{-}X}\NormalTok{ POST }\AttributeTok{{-}d} \StringTok{"name=john"}\NormalTok{ https://api.example.com/users  }\CommentTok{\# POST}
\ExtensionTok{curl} \AttributeTok{{-}C} \AttributeTok{{-}} \AttributeTok{{-}o}\NormalTok{ large.iso https://example.com/large.iso       }\CommentTok{\# Resume download}
\ExtensionTok{curl} \AttributeTok{{-}s}\NormalTok{ https://api.ipify.org                    }\CommentTok{\# Silent mode}
\ExtensionTok{curl} \AttributeTok{{-}I}\NormalTok{ https://example.com                      }\CommentTok{\# Check headers}
\end{Highlighting}
\end{Shaded}

\textbf{wget --- The Download Specialist:}

\begin{Shaded}
\begin{Highlighting}[]
\FunctionTok{wget}\NormalTok{ https://example.com/file.zip                }\CommentTok{\# Simple download}
\FunctionTok{wget} \AttributeTok{{-}P}\NormalTok{ \textasciitilde{}/Downloads https://example.com/file.zip }\CommentTok{\# To specific dir}
\FunctionTok{wget} \AttributeTok{{-}c}\NormalTok{ https://example.com/large.iso            }\CommentTok{\# Resume}
\FunctionTok{wget} \AttributeTok{{-}b}\NormalTok{ https://example.com/large{-}file.tar.gz    }\CommentTok{\# Background}
\FunctionTok{wget} \AttributeTok{{-}{-}limit{-}rate}\OperatorTok{=}\NormalTok{500k https://example.com/file.zip }\CommentTok{\# Limit speed}
\FunctionTok{wget} \AttributeTok{{-}r} \AttributeTok{{-}A.pdf}\NormalTok{ https://example.com/documents/    }\CommentTok{\# Recursive download}
\end{Highlighting}
\end{Shaded}

\textbf{aria2 --- The Download Accelerator:}

\begin{Shaded}
\begin{Highlighting}[]
\ExtensionTok{aria2c} \AttributeTok{{-}x}\NormalTok{ 16 }\AttributeTok{{-}s}\NormalTok{ 16 https://example.com/large.iso }\CommentTok{\# Multi{-}connection (faster!)}
\ExtensionTok{aria2c} \AttributeTok{{-}c}\NormalTok{ https://example.com/large.iso           }\CommentTok{\# Resume}
\ExtensionTok{aria2c} \AttributeTok{{-}{-}max{-}download{-}limit}\OperatorTok{=}\NormalTok{2M https://example.com/file.zip}
\end{Highlighting}
\end{Shaded}

\begin{quote}
\textbf{curl vs wget:} Use \texttt{curl} for APIs and complex HTTP. Use
\texttt{wget} for reliable downloads and recursive fetching.
\end{quote}

\begin{center}\rule{0.5\linewidth}{0.5pt}\end{center}

\paragraph{4) Remote Access}\label{remote-access}

\textbf{ssh --- Secure Shell:}

\begin{Shaded}
\begin{Highlighting}[]
\FunctionTok{ssh}\NormalTok{ username@remote{-}host                         }\CommentTok{\# Basic connection}
\FunctionTok{ssh} \AttributeTok{{-}p}\NormalTok{ 2222 username@remote{-}host                 }\CommentTok{\# Specific port}
\FunctionTok{ssh}\NormalTok{ user@host }\StringTok{"ls {-}la /var/log"}                  \CommentTok{\# Execute command remotely}
\FunctionTok{ssh} \AttributeTok{{-}L}\NormalTok{ 8080:localhost:80 user@host               }\CommentTok{\# Tunnel}
\FunctionTok{ssh} \AttributeTok{{-}D}\NormalTok{ 1080 user@host                            }\CommentTok{\# SOCKS proxy}
\FunctionTok{ssh} \AttributeTok{{-}X}\NormalTok{ user@host                                 }\CommentTok{\# X11 forwarding}
\ExtensionTok{ssh{-}copy{-}id}\NormalTok{ user@host                            }\CommentTok{\# Key{-}based auth (no password!)}
\FunctionTok{ssh} \AttributeTok{{-}i}\NormalTok{ \textasciitilde{}/.ssh/my\_key user@host                   }\CommentTok{\# Use specific key}
\FunctionTok{ssh}\NormalTok{ user@robot{-}host}
\end{Highlighting}
\end{Shaded}

\textbf{SSH Key Generation:}

\begin{Shaded}
\begin{Highlighting}[]
\FunctionTok{ssh{-}keygen} \AttributeTok{{-}t}\NormalTok{ ed25519}
\end{Highlighting}
\end{Shaded}

\begin{quote}
\textbf{Next Steps:} Disable password authentication in
\texttt{/etc/ssh/sshd\_config}. Use key pairs with passphrase. Change
default port. Use \texttt{fail2ban} to block brute force attempts.
\end{quote}

\textbf{scp --- Secure Copy:}

\begin{Shaded}
\begin{Highlighting}[]
\FunctionTok{scp}\NormalTok{ file.txt user@host:/home/user/               }\CommentTok{\# Local to remote}
\FunctionTok{scp}\NormalTok{ user@host:/var/log/syslog ./local{-}syslog     }\CommentTok{\# Remote to local}
\FunctionTok{scp} \AttributeTok{{-}r}\NormalTok{ ./project user@host:/var/www/             }\CommentTok{\# Directory}
\FunctionTok{scp} \AttributeTok{{-}p}\NormalTok{ file.txt user@host:/backup/               }\CommentTok{\# Preserve attributes}
\FunctionTok{scp}\NormalTok{ model.bin user@robot{-}host:/tmp/}
\end{Highlighting}
\end{Shaded}

\textbf{rsync --- Efficient Incremental Sync:}

\begin{Shaded}
\begin{Highlighting}[]
\FunctionTok{rsync} \AttributeTok{{-}av}\NormalTok{ /src/ /dest/                           }\CommentTok{\# Local sync}
\FunctionTok{rsync} \AttributeTok{{-}avz}\NormalTok{ /src/ user@host:/dest/                }\CommentTok{\# Remote sync}
\FunctionTok{rsync} \AttributeTok{{-}av}\NormalTok{ logs/ user@robot{-}host:/tmp/logs/}
\end{Highlighting}
\end{Shaded}

\begin{quote}
Learn the difference trailing slash makes---it's the ``slash of power''
vs.
\end{quote}

\begin{center}\rule{0.5\linewidth}{0.5pt}\end{center}

\paragraph{5) Multi-Robot Network
Configuration}\label{multi-robot-network-configuration}

Plan IP ranges, naming conventions, and role-based segmentation.

\begin{center}\rule{0.5\linewidth}{0.5pt}\end{center}

\paragraph{6) Static IP Setup for Raspberry Pi
Robots}\label{static-ip-setup-for-raspberry-pi-robots}

Stable addressing is critical for predictable remote control and
orchestration.

\begin{center}\rule{0.5\linewidth}{0.5pt}\end{center}

\paragraph{7) SSH Key-Based
Authentication}\label{ssh-key-based-authentication}

Use keys for secure, passwordless automation across robot fleets.

\begin{center}\rule{0.5\linewidth}{0.5pt}\end{center}

\paragraph{Useful Links --- Session 10}\label{useful-links-session-10}

\textbf{Bash Scripting:}

%
  \par\vspace{0.5ex}%
  \href{https://youtu.be/tK9Oc6AEnR4?si=y3M0s1AIVfqtzCJB}{\textcolor{blue}{\underline{\faLink\quad Bash scripting tutorial (YouTube)}}}%
  \par\vspace{0.5ex}%

%
  \par\vspace{0.5ex}%
  \href{https://mywiki.wooledge.org/BashGuide}{\textcolor{blue}{\underline{\faLink\quad BashGuide}}}%
  \par\vspace{0.5ex}%

%
  \par\vspace{0.5ex}%
  \href{https://www.shellcheck.net/}{\textcolor{blue}{\underline{\faLink\quad ShellCheck — lint your scripts}}}%
  \par\vspace{0.5ex}%

%
  \par\vspace{0.5ex}%
  \href{https://explainshell.com/}{\textcolor{blue}{\underline{\faLink\quad ExplainShell — understand complex commands}}}%
  \par\vspace{0.5ex}%

\textbf{Networking Tools:}

%
  \par\vspace{0.5ex}%
  \href{https://curl.se/docs/}{\textcolor{blue}{\underline{\faLink\quad curl documentation}}}%
  \par\vspace{0.5ex}%

%
  \par\vspace{0.5ex}%
  \href{https://www.gnu.org/software/wget/manual/wget.html}{\textcolor{blue}{\underline{\faLink\quad GNU Wget Manual}}}%
  \par\vspace{0.5ex}%

%
  \par\vspace{0.5ex}%
  \href{https://aria2.github.io/manual/en/html/}{\textcolor{blue}{\underline{\faLink\quad aria2 manual}}}%
  \par\vspace{0.5ex}%

%
  \par\vspace{0.5ex}%
  \href{https://www.openssh.com/manual.html}{\textcolor{blue}{\underline{\faLink\quad OpenSSH Manual}}}%
  \par\vspace{0.5ex}%

\textbf{Man Pages:}

%
  \par\vspace{0.5ex}%
  \href{https://linux.die.net/man/8/ping}{\textcolor{blue}{\underline{\faLink\quad ping}}}%
  \par\vspace{0.5ex}%

%
  \par\vspace{0.5ex}%
  \href{https://man7.org/linux/man-pages/man8/ip.8.html}{\textcolor{blue}{\underline{\faLink\quad ip}}}%
  \par\vspace{0.5ex}%

%
  \par\vspace{0.5ex}%
  \href{https://linux.die.net/man/8/ifconfig}{\textcolor{blue}{\underline{\faLink\quad ifconfig}}}%
  \par\vspace{0.5ex}%

%
  \par\vspace{0.5ex}%
  \href{https://man7.org/linux/man-pages/man1/scp.1.html}{\textcolor{blue}{\underline{\faLink\quad scp}}}%
  \par\vspace{0.5ex}%

\begin{center}\rule{0.5\linewidth}{0.5pt}\end{center}

\subsubsection{Session 11: Building Embedded Linux Systems with
Buildroot}\label{session-11-building-embedded-linux-systems-with-buildroot}

\begin{quote}
\textbf{Requirements:} 20 GB+ free storage, stable internet\\
\textbf{Tutorial reference:}
https://youtu.be/ey3sKdOmPa8?si=H0kdd-9R5Id4XoXy
\end{quote}

\paragraph{Learning outcomes}\label{learning-outcomes-10}

\begin{itemize}
\tightlist
\item
  Explain Buildroot vs Yocto tradeoffs
\item
  Configure and build a minimal embedded Linux image
\item
  Interpret output artifacts (kernel/rootfs/bootloader)
\item
  Apply basic customization and package integration
\end{itemize}

\begin{center}\rule{0.5\linewidth}{0.5pt}\end{center}

\paragraph{1) Build Systems Overview}\label{build-systems-overview}

\textbf{Buildroot:} simpler, faster, approachable for learning and many
practical projects\\
\textbf{Yocto:} highly flexible and powerful, but steeper complexity

\begin{center}\rule{0.5\linewidth}{0.5pt}\end{center}

\paragraph{2) Buildroot Architecture and
Workflow}\label{buildroot-architecture-and-workflow}

Pipeline: 1. Get source 2. Select base configuration 3. Configure
options 4. Build toolchain + packages + images 5. Collect final
artifacts

\begin{center}\rule{0.5\linewidth}{0.5pt}\end{center}

\paragraph{3) Configuration Basics}\label{configuration-basics}

Start with a baseline, then customize:

\begin{Shaded}
\begin{Highlighting}[]
\FunctionTok{make}\NormalTok{ defconfig          }\CommentTok{\# Baseline setup}
\FunctionTok{make}\NormalTok{ menuconfig         }\CommentTok{\# Interactive customization}
\end{Highlighting}
\end{Shaded}

\textbf{Core choices:} target architecture, toolchain, init system,
filesystem image format.

\begin{center}\rule{0.5\linewidth}{0.5pt}\end{center}

\paragraph{4) First Build and Outputs}\label{first-build-and-outputs}

Run the build and inspect the output:

\begin{Shaded}
\begin{Highlighting}[]
\FunctionTok{make}                    \CommentTok{\# Build everything}
\FunctionTok{ls}\NormalTok{ output/images        }\CommentTok{\# Check results}
\end{Highlighting}
\end{Shaded}

\texttt{make} produces images such as: - Root filesystem - Kernel image
- Bootloader-related artifacts

\begin{center}\rule{0.5\linewidth}{0.5pt}\end{center}

\paragraph{5) Customization}\label{customization}

\begin{itemize}
\tightlist
\item
  Add packages through menuconfig
\item
  Use rootfs overlays for custom files
\item
  Add post-build scripts
\item
  Create custom packages with \texttt{.mk} and \texttt{Config.in}
\end{itemize}

\begin{center}\rule{0.5\linewidth}{0.5pt}\end{center}

\paragraph{6) Practical Exercise}\label{practical-exercise}

Build a minimal ARM system for Raspberry Pi emulation including Python
and SSH support.

\begin{Shaded}
\begin{Highlighting}[]
\FunctionTok{make}\NormalTok{ defconfig}
\FunctionTok{make}\NormalTok{ menuconfig}
\FunctionTok{make}
\FunctionTok{ls}\NormalTok{ output/images}
\end{Highlighting}
\end{Shaded}

\begin{center}\rule{0.5\linewidth}{0.5pt}\end{center}

\paragraph{Useful Links --- Session 11}\label{useful-links-session-11}

\textbf{Buildroot:}

%
  \par\vspace{0.5ex}%
  \href{https://youtu.be/ey3sKdOmPa8?si=H0kdd-9R5Id4XoXy}{\textcolor{blue}{\underline{\faLink\quad Buildroot tutorial video}}}%
  \par\vspace{0.5ex}%

\textbf{Arch Wiki:}

%
  \par\vspace{0.5ex}%
  \href{https://wiki.archlinux.org/}{\textcolor{blue}{\underline{\faLink\quad Arch Wiki — the ultimate Linux reference}}}%
  \par\vspace{0.5ex}%

\begin{center}\rule{0.5\linewidth}{0.5pt}\end{center}

\paragraph{Useful Links --- Cool Stuff \&
Extras}\label{useful-links-cool-stuff-extras}

\textbf{Multimedia:}

%
  \par\vspace{0.5ex}%
  \href{https://itsfoss.com/ffmpeg/}{\textcolor{blue}{\underline{\faLink\quad FFmpeg tutorial (It's FOSS)}}}%
  \par\vspace{0.5ex}%

%
  \par\vspace{0.5ex}%
  \href{https://ffmpeg.org/ffmpeg.html}{\textcolor{blue}{\underline{\faLink\quad FFmpeg docs}}}%
  \par\vspace{0.5ex}%

\textbf{FOSS Alternatives:}

%
  \par\vspace{0.5ex}%
  \href{https://www.neowin.net/guides/top-10-foss-apps-to-make-your-linux-experience-more-enjoyable/}{\textcolor{blue}{\underline{\faLink\quad Top 10 FOSS Apps (Neowin)}}}%
  \par\vspace{0.5ex}%

%
  \par\vspace{0.5ex}%
  \href{https://alternativeto.net/}{\textcolor{blue}{\underline{\faLink\quad AlternativeTo.net}}}%
  \par\vspace{0.5ex}%

\textbf{Windows Apps/Games on Linux:}

%
  \par\vspace{0.5ex}%
  \href{https://youtu.be/xGbQZt96q54?si=FgQJ\_\_YMpISs7L1c}{\textcolor{blue}{\underline{\faLink\quad Bottles + Wine (YouTube)}}}%
  \par\vspace{0.5ex}%

%
  \par\vspace{0.5ex}%
  \href{https://youtu.be/JXHvGgiJJG0?si=1fI81xIr7zichNDy}{\textcolor{blue}{\underline{\faLink\quad Lutris (YouTube)}}}%
  \par\vspace{0.5ex}%

\textbf{Sync \& Backup:}

%
  \par\vspace{0.5ex}%
  \href{https://www.digitalocean.com/community/tutorials/how-to-use-rsync-to-sync-local-and-remote-directories}{\textcolor{blue}{\underline{\faLink\quad DigitalOcean rsync tutorial}}}%
  \par\vspace{0.5ex}%

%
  \par\vspace{0.5ex}%
  \href{https://youtu.be/GqSxR93xK6E?si=c2TvpU8R-53SXDty}{\textcolor{blue}{\underline{\faLink\quad rsync backup tutorial (YouTube)}}}%
  \par\vspace{0.5ex}%

\textbf{Ricing \& Customization:}

%
  \par\vspace{0.5ex}%
  \href{https://youtu.be/z1-yuolwrVs?si=3IKioFlYz9L8vkNS}{\textcolor{blue}{\underline{\faLink\quad Ricing inspiration (YouTube)}}}%
  \par\vspace{0.5ex}%

%
  \par\vspace{0.5ex}%
  \href{https://github.com/mikeroyal/Perfect-Ubuntu-Guide}{\textcolor{blue}{\underline{\faLink\quad Perfect Ubuntu Guide (GitHub)}}}%
  \par\vspace{0.5ex}%

%
  \par\vspace{0.5ex}%
  \href{https://mmbesar.github.io/Tutorials/Best-Arabic-Font-in-Linux/}{\textcolor{blue}{\underline{\faLink\quad Best Arabic Font in Linux (mmbesar)}}}%
  \par\vspace{0.5ex}%

%
  \par\vspace{0.5ex}%
  \href{https://github.com/Tinram/Linux-Utilities}{\textcolor{blue}{\underline{\faLink\quad Linux Utilities (GitHub Tinram)}}}%
  \par\vspace{0.5ex}%

%
  \par\vspace{0.5ex}%
  \href{https://youtu.be/nCS4BtJ34-o?si=T6gagbT5k71H5BNE}{\textcolor{blue}{\underline{\faLink\quad 12 GREAT CLI programs (YouTube)}}}%
  \par\vspace{0.5ex}%

\textbf{Shell \& Scripting Resources:}

%
  \par\vspace{0.5ex}%
  \href{https://learnbyexample.github.io/scripting\_course/Linux\_curated\_resources.html}{\textcolor{blue}{\underline{\faLink\quad Shell scripting course (Learn By Example)}}}%
  \par\vspace{0.5ex}%

\begin{center}\rule{0.5\linewidth}{0.5pt}\end{center}

\subsection{Suggested Projects}\label{suggested-projects}

\begin{enumerate}
\def\labelenumi{\arabic{enumi}.}
\tightlist
\item
  \textbf{System Monitor Script:} Create a Bash monitor script for
  CPU/RAM with logging and \textgreater80\% alerts.
\item
  \textbf{Automated Backup:} Build an automated workspace backup script
  to remote/cloud storage.
\item
  \textbf{SSH Remote Execution:} Practice SSH remote command execution
  between two machines/VMs (or localhost).
\item
  \textbf{QEMU Emulation:} Run full Raspberry Pi OS emulation in QEMU.
\item
  \textbf{Mobile SSH:} Optional: manage your Linux machine remotely from
  mobile (Termux + SSH).
\item
  \textbf{Custom Buildroot Image:} Build and document a custom Buildroot
  image with Python + SSH.
\item
  \textbf{Info Script:} Write a bash script that fetches currency rates,
  prayer times, or weather forecast using APIs with colored output and
  error handling.
\item
  \textbf{Rice Your Desktop:} Create a customized desktop environment
  with documented dotfiles on Git.
\end{enumerate}

\begin{center}\rule{0.5\linewidth}{0.5pt}\end{center}

\subsection{Contact the Author}\label{contact-the-author}

\begin{itemize}
\tightlist
\item
  \textbf{Email:} khaledmhfz2004@gmail.com
\item
  \textbf{GitHub:} https://github.com/KhaledMahfouz5/linux-course
\end{itemize}

\begin{center}\rule{0.5\linewidth}{0.5pt}\end{center}

\subsection{Copyright}\label{copyright}

All rights reserved for \textbf{hamakRobotTeam} {[}NOT For Commercial
Use{]}!!

DONE !! thanks Allah.

\end{document}
